% Modernised reading — fonds Alexandre Grothendieck, shelfmark 133
% (pages 1-53, the whole folder).
%
% This is an INTERPRETATION, not a transcription. It re-reads the pages in
% today's notation and vocabulary, states the hypotheses the manuscript leaves
% implicit, and drops the critical apparatus so that the mathematics reads
% straight through. Everything the brackets and underlines carried moves into a
% footnote. It is held to being mathematically correct as it stands; where the
% manuscript is loose or mistaken, the true statement is what appears here and
% the footnote says what the page has. In any dispute of substance the
% transcription governs: whoever cites Grothendieck must cite batch-01.fr.tex
% (pp. 1-20), batch-02.fr.tex (pp. 21-40) or batch-03.fr.tex (pp. 41-53).
%
% Pass: Opus 5.5 (claude-opus-5-5), 2026-09-23 — first pass, unchecked against the pages by a human.
% The folder was read whole, its three batches together; this is one document
% for the shelfmark. It was made on Opus 5.5 and has not been measured against
% the reference reading of folder 115 (made on Opus 5). The literature named
% in the footnotes is cited from memory and was not checked against the
% sources.
%
% Scope. Pages 1, 34 and 44 are covers whose words head sections here. Pages
% 4, 9 and 40 are blank. Pages 10-16 are two printed offprints of O. Laborde
% (C. R. Acad. Sc. Paris 1974), with nothing in his hand; they are identified,
% not described. Page 17 is a computer listing carrying one formula of his.
% Pages 47 and 49 are typed notices of the USTL (49: an « avis de soutenance »
% dated 5 June 1975), versos of 46 and 48. The unity of the folder is the
% archivists' bracketed title, « [Groupes de Witt et formes quadratiques,
% groupes formels] »; it is not his, and the reading says so.
%
% Spine. Seven runs, unrelated except where noted; kept distinct.
%   pp. 1-3     cover « groupes Witt et formes quadratiques »; scratch: the
%               associativity square and Mac Lane's pentagon with an
%               obstruction in Ext^1(pi0,pi0,pi0;pi1); w1 = [-1], w2 = c1 mod 2
%               for the hyperbolic plane on L; the change-of-rings spectral
%               sequence.
%   pp. 5-8     « Groupes de Witt et variantes »: WG(S) for a locally ringed
%               topos, pre-lambda-ring, the lambda-ring question reduced to
%               B_{O(m)}; hyperbolic map alpha: K_s(S) -> WG(S), W(S) = coker;
%               « Laborde dixit »; total Stiefel-Whitney class. Offprints
%               pp. 10-16.
%   pp. 17-28   G-sets and symmetric groups: k_*(G) = sum k(G x S_n) with the
%               induction product; tau^n(E) = cl_n(Mon(I_n,E)), exponential;
%               T^n(E) = cl_n(E^{I_n}); functors from finite sets and
%               injections (special, distinguished, « equivalences »,
%               Omega = End); on a ringed topos, the tensor power in
%               K(X,*) = sum K(B_{S_n}X) as sum over types of m_alpha(lambda)
%               rho^alpha; lambda^n and Adams psi^n as coefficients.
%   pp. 29-33   the game of go: territory score vs area, D = 2(P+T) - nu^2 -
%               eps; rules of capture; end of the obligation to play; who
%               plays last.
%   pp. 34-39   cover « Groupes formels »; statements on p-divisible and
%               unipotent formal groups; the « Solution »: a group G seen
%               through Cent(G), D(G), two crossed modules sharing pi0, pi1;
%               scratch on central extensions, Ga-hat by Ga-hat.
%   pp. 41-43   smooth connected algebraic groups with a Chevalley
%               decomposition: G = L x^z Z, Z anti-affine; Ext^1(A, z) =
%               Hom(D(z), Pic^0); char 0 vs char p; non-constructibility of
%               G_aff; « Reformulation ».
%   pp. 44-53   cover « Fourbis auxiliaires cohomologiques pour groupes
%               formels », his pp. 1-7: realising two crossed modules by one
%               G; obstruction H^3, indeterminacy H^2; N' = Cent G via Psi_G,
%               N = DG via lambda_G; dependence on the central extension E
%               through its commutator pairing c_E.
%
% Notation fixed for the document.
%  - « extension de A par B » is Grothendieck's: 1 -> B -> E -> A -> 1.
%  - [-1] := partial(-1) in H^1(S, mu_2) (p. 3 writes alpha, which pp. 6-7 use
%    for the hyperbolic map). alpha(E) = H(E) = E + E-check.
%  - G-sets: k_0(G), k_n(G) = k(G x S_n), cl_n, tau^n (injections, p. 19),
%    T^n (all maps, p. 21). On pp. 27-28 the tensor power, which the page
%    calls tau^n, is written T^n: it linearises E^{I_n}, not Mon(I_n, E).
%  - Formal groups: Cent(G), D(G) (the page also writes Z(G) in gothic, Der G,
%    DG); the subgroup of M on pp. 46-51 is gothic Z; the group L cap Z of
%    pp. 41-43 is fraktur z; on p. 45 N cap N' is written pi_1 only. Z on
%    pp. 41-43 is the anti-affine subgroup, nothing else.
%  - Crossed module (N, M, d, Theta); pi_1 = Ker d, pi_0 = Coker d.
%
% Substantive departures, each footnoted where it occurs:
%  p. 2    the target « L » of the square read as L_.; the Ext groups read as
%          multilinear extension groups (ours).
%  p. 3    abutment Ext_A^{p+q}(M,N) supplied; 2 invertible supplied; w2 = c1
%          mod 2 is his claim with a « ? », left unproved.
%  p. 5    P_i takes lambda^1..lambda^i (page: ..lambda^{i-1}); even rank needs
%          a point of residue characteristic 2; the product without 2
%          invertible checked (ours).
%  p. 6    alpha does NOT factor through K(S) in general: Atiyah's F_2 on an
%          elliptic curve (ours).
%  p. 7    alpha(K(S)) read alpha(K_s(S)); Laborde's statement proved via a
%          bilinear lift of q (ours); the B_{O(n)} counterexample is his,
%          unproved; the lambda-structure does not pass to W(S) (ours).
%  p. 8    the struck « w o alpha = 1 » is false: w(H(O)) = 1 + [-1];
%          normalisation w = 1 + class supplied.
%  p. 18   sigma_{-t} lambda_t = 1, E^2 = sigma^2 + lambda^2 and 2 lambda^2 +
%          E = E^2 hold for cardinalities, fail in A(Z/2) (ours); only
%          sigma^2 = lambda^2 + E holds as G-sets. A(G) -> R(G) injective iff
%          G cyclic (the page draws a hook).
%  p. 21   the cancelled formula replaced by the correct one over types alpha
%          with the stabilisers W_alpha = prod S_i wr S_{alpha_i} (ours).
%  p. 24   special = direct summands of finite SUMS of H_I (page: of an H_I;
%          2 h_{1} is a counterexample).
%  p. 28   NB (linear independence of the rho^alpha) true over Q, false over
%          F_2 (rho_1^2 = 2 rho_2); Question: yes for fields of char 0, no for
%          F_p (ours); sum index i >= 0.
%  p. 29   « nu^2 - D' » read nu^2 + D'; hypotheses: no seki, no neutral
%          point, every dead stone captured.
%  p. 33   the rule presupposes that each side only fills its own territory.
%  pp. 35-36  the page's labels pair Z(G) with G/Z(G) and D(G) with G/D(G);
%          exactness forces Cent(G) -> G/D(G) and D(G) -> G/Cent(G) (p. 36's
%          arcs agree).
%  p. 41   Ext^1(A, z) = Hom(D(z), Pic^0_{A/k}) (page: Pic_{A/k}); k perfect.
%  p. 42   phi_0 : t-dual -> H^1(A, O_A) (page: t -> ...).
%  p. 43   over the algebraic closure of a finite field z_m must be finite too
%          (ours); the non-constructibility sharpened (ours).
%  p. 45   [N, N'] = 1 supplied (needed for the crossed-module structures).
%  p. 50   Psi_0 injective <=> N' = H cap Cent(G) (page: N' = Cent(H); only
%          one direction holds).
%  p. 53   the « vérifier » (lambda_0 onto <=> dN = DM) checked (ours).
%
% Links the pages do not write (ours, and said so in the body): p. 3 computes
% w(alpha(L)), which is why p. 8's « w o alpha = 1 » is struck; p. 2 and pp.
% 45-53 are the pi_0/pi_1/class-one-degree-up pattern of folder 135; p. 21's
% decomposition Hom(I_n,E) = sum Mon(Q,E) is p. 24's H_J = sum h_Q; pp. 27-28
% are the linear counterpart of pp. 21-22 and take up folder 2 (pp. 15-18);
% p. 8's target ring is folder 2's A-tilde; p. 35's « Solution » is the problem
% of pp. 45-53, and p. 36's phi_G is p. 46's; pp. 41-43 are the algebraic
% model of p. 38's scratch.
%
% Announced and never established: the meaning of p. 2's obstruction; w2 = c1
% (p. 3); whether WG(S) is a lambda-ring and the task WG(B_{O(m)}) (p. 5); the
% B_{O(n)} counterexample (p. 7); existence of w (p. 8); lambda^i(xy) on G-sets
% (p. 18); any use of q_H^n (p. 20); the relation T*/tau* (p. 21, cancelled);
% the permanence of distinguished functors (p. 25); what Omega is for (p. 26);
% the Ko rule « cf. plus bas » (p. 31); whether eps depends on the play
% (p. 30); every statement of p. 35 and its « Question »; p. 36's « G
% unipotent »; p. 37's « voir fin »; the « bien connu » of p. 42 and the NB of
% p. 43; the obstruction and indeterminacy of p. 45; « on doit trouver »
% (p. 48); the unfinished sentence of p. 51; which E satisfy c) (p. 53); any
% return to formal groups.
\documentclass[11pt,a4paper]{article}
% Le préambule commun à toutes les transcriptions du fonds Grothendieck.
%
% Il tient en une idée : **l'incertitude se compose, elle ne se résout pas.**
% Une transcription qui lisse un mot douteux a détruit la seule chose pour
% laquelle on la faisait — un lecteur doit pouvoir savoir, à chaque endroit,
% ce qui était sur le feuillet et ce qui a été ajouté. Les six macros qui
% suivent sont donc l'appareil critique, et non de la décoration.
%
% Les mêmes commandes sont reconnues par `scripts/render.mjs`, qui produit la
% vue de lecture du site. Une macro ajoutée ici doit l'être aussi là-bas,
% faute de quoi le rendu échouera bruyamment — ce qui est voulu : un rendu
% silencieusement faux est pire qu'un rendu absent.

\ProvidesPackage{grothendieck}[2026/08/08 Transcriptions du fonds Grothendieck]

\usepackage{amsmath,amssymb,amsthm}
\usepackage{mathrsfs}
\usepackage{stmaryrd}
% Les diagrammes commutatifs sont partout dans ce fonds : c'est la langue
% même de la géométrie algébrique de Grothendieck, et une transcription qui
% les rendrait en prose ne transcrirait rien.
\usepackage{tikz-cd}
% Français par défaut — c'est la langue des feuillets — mais l'anglais est
% chargé aussi : la traduction et le résumé sont des documents anglais, et la
% typographie française (espaces avant les deux-points) y serait une faute.
% Ces documents basculent par \selectlanguage{english} en tête de corps.
\usepackage[english,french]{babel}
\usepackage{fontspec}
\usepackage{geometry}
\geometry{a4paper,margin=25mm,marginparwidth=28mm}
\usepackage{marginnote}
\usepackage{xcolor}
% ulem, not soul. `soul`'s \st and \ul are built on a character-by-character
% scan that XeLaTeX's font machinery breaks: the document fails with an
% undefined control sequence rather than degrading. ulem does the same two jobs
% and is Unicode-engine safe. `normalem` stops it hijacking \emph, which the
% transcriptions use for Grothendieck's own emphasis.
\usepackage[normalem]{ulem}
\usepackage{hyperref}
\hypersetup{colorlinks=true,linkcolor=black,urlcolor=black,pdfborder={0 0 0}}

% --- Métadonnées du lot -----------------------------------------------------
% Reprises telles quelles par le script de rendu, qui les affiche en tête de
% la vue de lecture. Elles fixent ce que le fichier prétend couvrir : sans
% elles, une transcription flotte sans point d'attache dans un fonds de seize
% mille feuillets.

\newcommand{\folder}[1]{\gdef\grothFolder{#1}}
\newcommand{\batch}[1]{\gdef\grothBatch{#1}}
\newcommand{\leaves}[2]{\gdef\grothFirst{#1}\gdef\grothLast{#2}}
\newcommand{\foldertitle}[1]{\gdef\grothFoldertitle{#1}}
\gdef\grothFolder{?}\gdef\grothBatch{?}\gdef\grothFirst{?}\gdef\grothLast{?}\gdef\grothFoldertitle{}

% --- Le résumé --------------------------------------------------------------
%
% Il ouvre la lecture modernisée, sous le titre « Résumé », et il est composé
% en petit. Ce n'est pas un ornement : c'est le seul passage du document qui
% ne soit pas des mathématiques, et un lecteur doit pouvoir le reconnaître
% avant de l'avoir lu — puis le sauter s'il connaît déjà le sujet, ou s'y
% arrêter s'il ne le connaît pas. Le filet qui le suit marque la reprise du
% corps.

\newenvironment{resume}
  {\small\setlength{\parskip}{0.45em}}
  {\par\medskip\hrule\medskip}

% --- Mots-clés ---------------------------------------------------------------
%
% En anglais, en clôture du résumé de la lecture modernisée : le vocabulaire
% moderne sous lequel chercher ce que les feuillets construisent. Le manifeste
% du site (`npm run manifest`) les extrait pour étiqueter la cote entière —
% cette ligne est la source unique des tags, il n'y a pas de fichier à côté.
\newcommand{\keywords}[1]{\par\smallskip\noindent
  {\footnotesize\textsc{Keywords} --- \textit{#1}}\par}

% --- L'appareil critique ----------------------------------------------------

% Le numéro de feuillet, en marge. C'est l'ancre qui permet de régler un
% désaccord sur une formule en regardant *un* feuillet plutôt qu'une chemise
% de six cents. Le site s'en sert aussi pour tourner les pages du fac-similé
% quand on fait défiler la transcription.
\newcommand{\leaf}[1]{%
  \marginnote{\footnotesize\color{gray!70}\textsf{#1}}%
  \label{leaf:#1}\ignorespaces}

% Illisible : on ne devine pas. Un mot inventé qui se lit comme les autres est
% le pire résultat possible.
\newcommand{\ill}{\textcolor{red!60!black}{[\,\ldots\,]}}

% Lecture douteuse : le mot est donné, et signalé comme incertain.
\newcommand{\uncertain}[1]{\uline{#1}}

% Ajout de l'éditeur — une lettre manquante, un mot sous-entendu.
\newcommand{\add}[1]{\textcolor{blue!50!black}{[#1]}}

% Biffé par Grothendieck lui-même. Ce qu'il a rayé fait partie du document :
% une rature dit souvent où la pensée a changé de direction.
\newcommand{\struck}[1]{\sout{#1}}

% Note du transcripteur, sur l'état matériel du feuillet.
\newcommand{\note}[1]{\par\smallskip{\small\color{gray!60!black}\textit{[#1]}}\par\smallskip}

% Note marginale de Grothendieck, distincte du corps du texte.
\newcommand{\marginal}[1]{\par\smallskip\noindent\hspace{1em}%
  {\small\color{gray!60!black}#1}\par\smallskip}

% --- Titre ------------------------------------------------------------------

\AtBeginDocument{%
  \begin{center}
    {\large\bfseries Fonds Alexandre Grothendieck --- cote n\textsuperscript{o}~\grothFolder}\\[2pt]
    {\itshape\grothFoldertitle}\\[4pt]
    {\small feuillets \grothFirst--\grothLast\ (lot \grothBatch)}\\[2pt]
    {\footnotesize\color{gray!60!black}Transcription établie d'après le fac-similé numérisé
     par l'Université de Montpellier.\\
     Les lectures incertaines sont soulignées, les illisibles notées [\,\ldots\,],
     les ajouts entre crochets.}
  \end{center}
  \vspace{1em}}

\endinput


\folder{133}
\pages{1}{53}
\dating{1974-[à partir de 1975]}
\watermark{Édition de démonstration\\interprétation personnelle de l'œuvre}
\foldertitle{[Groupes de Witt et formes quadratiques, groupes formels] : notes manuscrites (s.d.), tirés à part (1974) — lecture modernisée du dossier entier}

\begin{document}

\section*{Résumé}

\begin{resume}
Ce dossier n'est pas un texte mais un paquet. Sept suites de notes, sans
date et sans lien entre elles, y sont rangées sous un titre que les
archivistes ont écrit entre crochets, « Groupes de Witt et formes
quadratiques, groupes formels » : c'est leur titre, non le sien. Trois
chemises de sa main découpent le paquet — « groupes Witt et formes
quadratiques », « Groupes formels », « Fourbis auxiliaires cohomologiques
pour groupes formels » — et entre elles se glissent une réflexion sur les
ensembles munis de symétries et quatre pages sur le jeu de go. On les lit ici
une par une, sans leur prêter une unité qu'elles n'ont pas.

La première suite part d'une question ancienne : comment classer les formes
quadratiques, ces expressions comme $x^{2} + y^{2}$ ou $xy$, non plus sur un
corps mais sur un espace qui varie. Grothendieck fait ce qu'on fait pour
passer des entiers naturels aux entiers relatifs : on autorise les
différences formelles de formes, puis on décide qu'une forme et son opposée
s'annulent. Ce qui reste est le groupe de Witt. Deux questions l'occupent.
Les opérations qui généralisent les coefficients binomiaux (les « puissances
extérieures ») obéissent-elles aux identités qu'elles vérifient pour les
nombres ? Il ne le sait pas, et réduit la question à un seul cas universel.
Et peut-on attacher à une forme des invariants de parité, les classes de
Stiefel–Whitney, qui mesurent en cohomologie modulo $2$ ce qui l'empêche
d'être une somme de carrés ?

La deuxième suite compte avec des symétries. Un ensemble fini sur lequel un
groupe agit a une « taille » plus fine qu'un nombre : l'ensemble des types
d'orbites qui le composent. Si l'on choisit $n$ éléments distincts d'un tel
ensemble, en les numérotant, on obtient un nouvel ensemble, sur lequel
agissent en plus les permutations des numéros ; Grothendieck montre que
cette construction transforme une réunion en un produit, comme
l'exponentielle transforme une somme en un produit. Il remarque au passage
que la vieille identité « $n^{2}$ couples $=$ $n$ couples diagonaux $+$ deux
fois les paires » cesse d'être vraie quand on retient la symétrie, et c'est
ce qui oblige à garder les permutations. Viennent ensuite les foncteurs
définis sur les ensembles finis et leurs injections, objets qu'on a
étudiés pour eux-mêmes une quarantaine d'années plus tard, puis la version linéaire : la puissance
tensorielle d'un module, avec les permutations de ses facteurs, contient à
la fois les puissances extérieures et les opérations d'Adams.

La suite sur le jeu de go est un exercice de comptabilité. Deux manières de
compter les points d'une partie — les prisonniers plus le territoire, ou les
pierres plus le territoire — donnent des résultats qui diffèrent d'au plus
un point, selon qui a joué le dernier ; Grothendieck l'établit, puis cherche
à dire précisément quand une partie est finie.

Les deux dernières suites, les plus longues, veulent reconstituer un groupe
non commutatif à partir de deux de ses sous-groupes : son centre et le
sous-groupe engendré par ses commutateurs. Chacun des deux fournit une
flèche dont le noyau et le conoyau sont les mêmes deux groupes, et la
question est de savoir si l'on peut recoller les deux flèches en un groupe.
L'analogie est celle de la retenue dans une addition à plusieurs chiffres :
on connaît chaque chiffre, et il faut encore choisir les retenues, de façon
cohérente. La réponse a la forme classique : un obstacle à l'existence,
dans un groupe de cohomologie de degré $3$, puis une liberté de choix,
mesurée en degré $2$ par les extensions centrales ; et il faut enfin que le
groupe construit ait \emph{exactement} le centre et les commutateurs
prescrits, ce que le dossier traduit en conditions sur une forme bilinéaire
alternée. Entre les deux, trois pages sur les groupes algébriques
décomposent un groupe en une partie affine et une variété abélienne, et
relèvent un phénomène qui sépare la caractéristique nulle de la
caractéristique $p$.

On voit ici des notes de travail au sens strict. Les ignorances sont écrites
comme telles (« j'ignore si c'est un $\lambda$-anneau », « il n'est pas
clair »), les affirmations d'autrui sont signées (« Laborde dixit »), et une
présentation exacte mais laide est rejetée pour une autre, parce qu'elle
oblige à manipuler des choses « trop dégueulasses » : le goût pour la donnée
naturelle plutôt que pour la donnée explicite. La datation vient de
l'inventaire ; les tirés à part sont de 1974, et un avis administratif de
juin 1975 sert de verso à l'une des dernières pages.

Les noms sous lesquels chercher la suite : groupe de Grothendieck–Witt et
groupe de Witt d'un schéma, classes de Stiefel–Whitney d'une forme
quadratique, $\lambda$-anneaux ; anneau de Burnside, opérations de
puissance, FI-modules et espèces de structures, dualité de Schur–Weyl ;
décompte par territoire et par zone au go ; modules croisés, obstruction
d'Eilenberg–Mac Lane, groupes anti-affines et formule de Barsotti–Weil.

\keywords{Grothendieck-Witt ring, Witt group of a scheme, nonsingular quadratic form, hyperbolic form, Stiefel-Whitney class of a quadratic form, lambda-ring, Mac Lane pentagon, change-of-rings spectral sequence, Burnside ring, power operation, FI-module, combinatorial species, monomial symmetric function, Schur-Weyl duality, Adams operation, rules of Go, area scoring, territory scoring, formal group, p-divisible group, crossed module, Eilenberg-Mac Lane obstruction, extension with non-abelian kernel, commutator pairing, Chevalley structure theorem, anti-affine group, Cartier duality, Barsotti-Weil formula}
\end{resume}

\pagerange{1}{53}
\section*{Le fil du dossier, et les conventions}

\subsection*{Les stations}

L'inventaire date le dossier « 1974-[à partir de 1975] » et lui donne un
titre entre crochets, qui est le leur\footnote{Le titre de l'inventaire est
« [Groupes de Witt et formes quadratiques, groupes formels] : notes
manuscrites (s.d.), tirés à part (1974) ». Les crochets sont ceux des
archivistes : les mots viennent vraisemblablement des deux premières
chemises, pages 1 et 34, mais la réunion des deux est de leur fait. Aucune
page ne relie les formes quadratiques aux groupes formels.}. Les pages
transcrites se répartissent en sept suites.

\begin{enumerate}
  \item \emph{Pages 1 à 3.} La chemise « groupes Witt et formes
  quadratiques », puis deux brouillons : l'associativité d'un produit sur un
  complexe, une obstruction et le pentagone de Mac Lane (page 2) ; les
  classes de Stiefel–Whitney du plan hyperbolique sur un fibré en droites,
  et la suite spectrale de changement d'anneaux (page 3).
  \item \emph{Pages 5 à 8}, « Groupes de Witt et variantes » : le groupe
  $WG(S)$ des modules quadratiques sur un topos localement annelé, la
  question de sa structure de $\lambda$-anneau, l'application hyperbolique
  et le groupe de Witt $W(S)$, la classe de Stiefel–Whitney totale. Les
  pages 10 à 16 sont deux tirés à part d'O.~Laborde, que la page 7 cite.
  \item \emph{Pages 17 à 28}, sans titre : l'anneau $k_{*}(G)$ des
  $G \times \mathfrak{S}_n$-ensembles finis virtuels, les opérations
  $\tau^{n}$ et $T^{n}$, les foncteurs sur la catégorie des ensembles finis
  et injections, puis les puissances tensorielles sur un topos annelé et les
  représentations des groupes symétriques.
  \item \emph{Pages 29 à 33}, le jeu de go : décompte des points, règles de
  prise, fin de l'obligation de jouer.
  \item \emph{Pages 34 à 39}, « Groupes formels » : énoncés sur les groupes
  formels $p$-divisibles et unipotents, et la « Solution » qui voit un groupe
  à travers son centre et son groupe dérivé ; brouillons.
  \item \emph{Pages 41 à 43}, sans titre : groupes algébriques lisses
  connexes admettant une décomposition de Chevalley, et leurs extensions
  centrales.
  \item \emph{Pages 44 à 53}, « Fourbis auxiliaires cohomologiques pour
  groupes formels », paginées par lui de 1 à 7 : réaliser deux modules
  croisés de mêmes invariants par un seul groupe, puis imposer que ce groupe
  ait le centre et le groupe dérivé voulus.
\end{enumerate}

Les pages 4, 9 et 40 sont blanches. La page 17 est un feuillet de listing
d'ordinateur qui ne porte de sa main qu'une formule. Les pages 47 et 49 sont
des avis dactylographiés de l'administration de l'université, versos des
pages 46 et 48 ; le second, un avis de soutenance daté du 5 juin 1975, lui
est adressé\footnote{Ils ne sont pas de sa main et ne sont pas transcrits.
Comme sa pagination court sans interruption de la page 46 (sa page 2) à la
page 48 (sa page 3) et à la page 50 (sa page 4), ces feuillets sont du papier
de réemploi ; l'avis de 1975 suggère seulement que la dernière suite n'est
pas antérieure à juin 1975.}.

\subsection*{Sept suites, gardées distinctes}

Les suites ne se citent pas : aucune page ne renvoie à une autre, et les
pages 41 à 43, sur les groupes algébriques, sont rangées entre les deux
suites sur les groupes formels sans qu'aucune ne les mentionne. Les liens
signalés plus bas sont donc tous de nous. Deux d'entre eux s'appuient
cependant sur ce que les pages montrent : la page 3 calcule la classe de
Stiefel–Whitney d'un plan hyperbolique, et c'est l'affirmation qu'une telle
classe est triviale qui est barrée page 8 ; la « Solution » de la page 35 et
les « Fourbis » des pages 45 à 53 donnent à la même application le même nom,
$\varphi_G$.

\subsection*{Conventions}

\begin{itemize}
  \item \emph{Extensions.} « Extension de $A$ par $B$ » a le sens de
  Grothendieck : une suite exacte $1 \to B \to E \to A \to 1$.
  \item \emph{Formes quadratiques.} Un module quadratique $(M, q)$ est un
  module localement libre de type fini muni d'une forme quadratique, de
  forme polaire $b(x, y) = q(x+y) - q(x) - q(y)$ ; il est
  \emph{strictement non dégénéré} (le mot de la page ; on dit aujourd'hui
  \emph{non singulier}) si $b$ identifie $M$ à son dual $\check{M}$. Le plan
  hyperbolique sur $E$ est $H(E) = E \oplus \check{E}$, $q(x + \xi) =
  \langle x, \xi \rangle$, et l'application hyperbolique est
  $\alpha(E) = H(E)$. La classe de $-1$ dans $H^{1}(S, \mu_2)$ est notée
  $[-1]$ ; la page 3 l'appelle $\alpha$, lettre que les pages 6 et 7
  réservent à l'application hyperbolique.
  \item \emph{Ensembles finis.} $I_n = \{1, \dots, n\}$ ; $\mathrm{Mon}(I, E)$
  est l'ensemble des injections, $E^{I} = \mathrm{Hom}(I, E)$ celui de
  toutes les applications. $\tau^{n}(E)$ est la classe de
  $\mathrm{Mon}(I_n, E)$ (page 19), $T^{n}(E)$ celle de $E^{I_n}$ (page 21).
  Aux pages 27 et 28, la puissance tensorielle $n$-ième d'un module, que la
  page appelle $\tau^{n}$, est notée ici $T^{n}$ : elle linéarise
  $E^{I_n}$ et non $\mathrm{Mon}(I_n, E)$, puisque l'algèbre d'un produit
  est le produit tensoriel des algèbres.
  \item \emph{Groupes.} $\mathrm{Cent}(G)$ est le centre, $D(G)$ le groupe
  dérivé (la page écrit aussi $\mathfrak{Z}(G)$, $\mathrm{Dér}\,G$, $DG$).
  Un module croisé est un quadruplet $(N, M, d, \Theta)$, d'invariants
  $\pi_1 = \operatorname{Ker} d$ et $\pi_0 = \operatorname{Coker} d$. Le
  sous-groupe de $M$ des pages 46 à 51 est noté $\mathfrak{Z}$ ; le groupe
  $L \cap Z$ des pages 41 à 43 est noté $\mathfrak{z}$, et $Z$ n'y désigne
  que le sous-groupe anti-affine. À la page 45, $N \cap N'$ est noté
  $\pi_1$ seulement\footnote{La page 45 l'écrit aussi avec le « 3 » bouclé
  qu'on rend par $\mathfrak{z}$ ; c'est un autre objet que le $\mathfrak{z}$
  des pages 41 à 43.}.
\end{itemize}

\subsection*{Ce que seule une lecture d'ensemble peut dire}

Six liens, qu'aucune page n'écrit, et qui sont de nous.

Le brouillon de la page 3 donne, pour le plan hyperbolique sur un fibré en
droites $L$, $w_1 = [-1]$ et $w_2 = c_1(L) \bmod 2$. C'est exactement
$w(\alpha(L))$, et cela montre que la phrase de la page 8 selon laquelle
$w \circ \alpha = 1$ est fausse : c'est celle que Grothendieck a barrée.

Le brouillon de la page 2 (des invariants $\pi_0$, $\pi_1$, une contrainte
d'associativité, une obstruction un degré plus haut modulo l'image d'un
degré plus bas) et les pages 45 à 53 (deux modules croisés d'invariants
$\pi_0$, $\pi_1$, une obstruction dans $H^{3}$, une indétermination dans
$H^{2}$) relèvent du même schéma : celui des Gr-catégories du dossier 135,
daté des mêmes années.

La décomposition $\mathrm{Hom}(I_n, E) = \coprod_{Q} \mathrm{Mon}(Q, E)$,
sur les quotients $Q$ de $I_n$, qui ouvre la page 21, est la décomposition
$H_J = \coprod_Q h_Q$ de la page 24, vue ensemble par ensemble ; c'est elle
qui fait des $T^{n}$ des polynômes en les $\tau^{m}$.

Les pages 27 et 28 sont la contrepartie linéaire des pages 21 et 22 : les
deux sommes y sont indexées par les types $\alpha$ de partitions de $n$. Et
elles reprennent le cadre du dossier 2 (pages 15 à 18), où le produit
d'induction sur les groupes symétriques et les puissances tensorielles
d'un module sur un topos annelé sont déjà posés.

La « Solution » de la page 35 et les « Fourbis » des pages 45 à 53 sont un
même problème : $\mathrm{Cent}(G)$ et $D(G)$ fournissent deux modules
croisés de mêmes invariants, et il s'agit de dire quand deux tels modules
croisés proviennent d'un groupe. Le point d'exclamation de la page 36
(« est quasi-unipotent ! ») est l'injectivité de $\Psi_G$ que la page 46
caractérise.

Les pages 41 à 43 donnent, pour les groupes algébriques, le modèle de ce que
le brouillon de la page 38 esquisse pour les groupes formels : une extension
centrale d'un groupe linéaire (ou de son complété formel) par un groupe
commutatif.

\subsection*{Ce que le dossier annonce sans l'établir}

Le sens de l'obstruction de la page 2. L'égalité $w_2 = c_1 \bmod 2$ de la
page 3, marquée d'un « ? ». La question de la page 5 (est-ce un
$\lambda$-anneau ?) et la tâche de déterminer $WG(B_{O(m)})$. Le
contre-exemple de la page 7 sur $B_{O(n)}$. L'existence de la classe $w$ de
la page 8. Les formules de $\lambda^{i}(xy)$ pour les $G$-ensembles (page
18). Tout usage des applications $q^{n}_H$ (page 20). La relation entre
$T^{*}$ et $\tau^{*}$ annoncée page 21, dont la formule est annulée. Les
propriétés de permanence des foncteurs distingués (page 25). Ce à quoi doit
servir $\Omega$ (page 26). La description du ko, « cf.\ plus
bas » (page 31), jamais écrite. La question de la page 30 : la valeur de
$\varepsilon$ dépend-elle de la suite du jeu ? Tous les énoncés de la page
35 et sa « Question ». Le « $G$ unipotent » de la page 36. Le « voir fin »
de la page 37, qu'aucune fin ne reprend. Le « bien connu » de la page 42 et
le NB de la page 43. L'obstruction et l'indétermination de la page 45. Le
« on doit trouver » de la page 48. La phrase inachevée de la page 51. La
détermination des extensions $E$ qui satisfont les conditions c) (pages 51
et 53). Et tout retour aux groupes formels : le dossier s'arrête sur la
page 53 sans y revenir.

\pagerange{1}{3}
\section*{I. « Groupes Witt et formes quadratiques » : la chemise et deux brouillons (pages 1 à 3)}

La page 1 est une chemise de carton rose qui ne porte que ces mots, en haut à
droite\footnote{Le dernier mot est abrégé, « qu.dr.tiques », et lu
« quadratiques ».}. Les deux feuillets qui suivent sont des brouillons, sans
phrase suivie.

\pagerange{2}{2}
\subsection*{Associativité et pentagone (page 2)}

Le feuillet, à en-tête de l'Institut de mathématiques de Montpellier, part
d'un complexe $L_{\bullet}$ muni d'un produit $p \colon L_{\bullet} \otimes
L_{\bullet} \to L_{\bullet}$, et dessine le carré qui exprime
l'associativité de $p$ : les deux composés
\[
p \circ (p \otimes L_{\bullet}) \quad \text{et} \quad
p \circ (L_{\bullet} \otimes p) \circ a,
\qquad
(L_{\bullet} \otimes L_{\bullet}) \otimes L_{\bullet} \longrightarrow
L_{\bullet},
\]
où $a$ est l'isomorphisme d'associativité du produit tensoriel des
complexes\footnote{La page écrit $L$ sans indice au sommet inférieur du
carré ; on lit $L_{\bullet}$, que le reste du carré demande. Une flèche
courbe entre les deux $L_{\bullet} \otimes L_{\bullet}$, marquée d'un
$\alpha$ douteux, est biffée.}. Sous le carré, une flèche
\[
\mathrm{Ext}(\pi_0, \pi_0 ; \pi_1) \xrightarrow{\ \delta^{(1)}\ }
\mathrm{Ext}^{1}(\pi_0, \pi_0, \pi_0 ; \pi_1).
\]
La page ne définit aucun de ces groupes. Notre lecture est la suivante, et
elle n'est qu'une lecture : si $\pi_0$ et $\pi_1$ sont les deux groupes
d'homologie d'un complexe à deux termes et si $p$ n'est associatif qu'à
homotopie près, le défaut d'associativité définit une classe dans un groupe
d'extensions « trilinéaires » de $\pi_0$ par $\pi_1$, et changer $p$ la
modifie par l'image, sous une différentielle $\delta^{(1)}$, d'un groupe
d'extensions « bilinéaires ». Seule la classe dans le conoyau de
$\delta^{(1)}$ mesure l'obstruction\footnote{C'est le schéma des
obstructions d'associativité qu'on rencontre en cohomologie de Hochschild
et de Mac Lane ; l'identification des groupes de la page est de nous. En
tête du feuillet, deux relations isolées, $q' + q'' + q''' = -1$ et
$q''' = q'q''$ (le signe central est surchargé), ne se rattachent à rien.}.

Le feuillet se termine sur le \emph{pentagone} : pour quatre objets
$x, y, z, t$, les deux chemins de $((x \otimes y) \otimes z) \otimes t$ à
$x \otimes (y \otimes (z \otimes t))$, l'un par $(x \otimes y) \otimes
(z \otimes t)$, l'autre par $(x \otimes (y \otimes z)) \otimes t$ et
$x \otimes ((y \otimes z) \otimes t)$, doivent coïncider :
\[
a_{x, y, z \otimes t} \circ a_{x \otimes y, z, t}
= (x \otimes a_{y, z, t}) \circ a_{x, y \otimes z, t} \circ
(a_{x, y, z} \otimes t).
\]
C'est l'axiome de cohérence de Mac Lane (1963), qui garantit que toutes les
manières de réassocier un produit de $n$ facteurs
s'accordent\footnote{La page dessine les cinq sommets sans écrire
l'égalité ; la flèche de droite, entre $(x \otimes y) \otimes (z \otimes t)$
et $x \otimes (y \otimes (z \otimes t))$, porte deux pointes. S.~Mac Lane,
« Natural associativity and commutativity », 1963 : Grothendieck ne le cite
pas ici, mais l'article est antérieur de dix ans.}. Le même schéma — deux
invariants, une contrainte, une classe un degré plus haut — est celui des
Gr-catégories du dossier 135 (pages 40 et 41), et on le retrouvera aux pages
45 à 53.

\pagerange{3}{3}
\subsection*{Une classe de Stiefel–Whitney, une suite spectrale (page 3)}

On suppose $2$ inversible\footnote{Hypothèse implicite de la page : elle
identifie $\mu_2$ à $\mathbb{F}_2$ et utilise la suite de Kummer
$1 \to \mu_2 \to \mathbb{G}_m \xrightarrow{2} \mathbb{G}_m \to 1$.}. Soit $L$
un fibré en droites sur $S$ et $E = L \oplus L^{-1}$, muni de la forme
quadratique $Q(x, x') = x x'$ : c'est le plan hyperbolique $H(L)$ des
pages 6 et 7. La page affirme (« Je dis que ») :
\[
w_1(E) = [-1] := \partial(-1), \qquad w_2(E) \equiv c_1(L) \bmod 2,
\]
où $\partial \colon H^{0}(S, \mathbb{G}_m) \to H^{1}(S, \mu_2)$ et
$H^{1}(S, \mathbb{G}_m) \to H^{2}(S, \mu_2)$ sont les cobords de Kummer. La
première égalité est juste : localement, $H(L)$ est la forme
$\langle 1, -1 \rangle$, de discriminant $-1$. La seconde est vérifiée
formellement : si $y$ est une « racine » de $E$, l'autre est $y + [-1]$, et
\[
(1 + y)(1 + y + [-1]) = 1 + [-1] + y(y + [-1]),
\]
la page inscrivant sous $y(y + [-1])$ « $c_1$ ? ». Nous ne l'établissons
pas ici\footnote{Deux vérifications sont compatibles avec l'énoncé : pour
$L$ trivial, $w_2(\langle 1, -1 \rangle) = 0 = c_1$ ; et sur les nombres
complexes, la réduction de $H(L)$ au groupe orthogonal compact est le fibré
réel orienté sous-jacent à $L$, dont la classe $w_2$ est $c_1 \bmod 2$. La
page écrit $\alpha$ pour $[-1]$.}.

En tête du feuillet, une classe $\eta \in H^{1}(P', \mathbb{F}_2)$ sur un
$P' = W^{*}(L^{\otimes 2})$ au-dessus duquel $L^{\otimes 2}$ devient trivial,
soumise à
\[
\eta^{2} + w_1 \eta + w_2 = 0 .
\]
La page ne dit pas ce qu'est $W^{*}$. La relation a la forme de celle par
laquelle son article de 1958 définit les classes de Chern sur le fibré
projectif, réduite modulo $2$ et en rang $2$\footnote{« La théorie des
classes de Chern », Bull. SMF 86 (1958) : sur le fibré projectif d'un fibré
de rang $r$, la classe tautologique $\xi$ vérifie
$\sum_{i} (-1)^{i} c_i \xi^{r-i} = 0$, et cette relation définit les $c_i$.
Le rapprochement est de nous.}.

Le bas du feuillet note, pour un homomorphisme d'anneaux $A \to B$, un
$A$-module $M$ et un $B$-module $N$, la suite spectrale de changement
d'anneaux
\[
E_2^{p,q} = \mathrm{Ext}_B^{p}\bigl(\mathrm{Tor}_q^{A}(M, B), N\bigr)
\Longrightarrow \mathrm{Ext}_A^{p+q}(M, N),
\]
issue de l'adjonction $\mathrm{Hom}_A(M, N) = \mathrm{Hom}_B(M \otimes_A B,
N)$\footnote{La page n'écrit que le terme $E_2$, les indices $p$ et $q$ en
surcharge sur un $B$ et un $A$ ; l'aboutissement est ajouté ici. C'est la
suite de Cartan et Eilenberg (\emph{Homological Algebra}, 1956). Quelques
lignes non mathématiques au pied du feuillet ne sont pas transcrites.}.

\pagerange{5}{16}
\section*{II. Groupes de Witt et variantes (pages 5 à 16)}

\pagerange{5}{5}
\subsection*{Le groupe $WG(S)$ et la question du $\lambda$-anneau (page 5)}

Soit $S$ un topos localement annelé. On note $WG(S)$ le groupe de
Grothendieck du monoïde, pour la somme orthogonale, des classes
d'isomorphie de modules quadratiques strictement non dégénérés sur $S$.
C'est ce qu'on appelle aujourd'hui le \emph{groupe de Grothendieck–Witt} de
$S$\footnote{Pour un anneau commutatif, M.~Knebusch introduit les
« Grothendieck- und Wittringe » en 1969-1970 ; pour un schéma, la
définition moderne (Knebusch, 1977, puis Walter et Schlichting) ajoute les
relations issues des sous-modules lagrangiens, que la définition de la page
ne contient pas. Hors du cas affine, les deux groupes peuvent différer ;
voir la note sur la page 6.}.

\emph{Structure.} Le produit tensoriel et les puissances extérieures font de
$WG(S)$ un pré-$\lambda$-anneau, augmenté par le rang vers l'anneau
binomial $H^{0}(S, \mathbb{Z})$ des fonctions localement constantes à
valeurs entières. En général, il n'a pas d'unité : l'unité naturelle serait
la forme $\langle 1 \rangle$, $q(x) = x^{2}$, dont la forme polaire $2xy$
n'est non dégénérée que si $2$ est inversible\footnote{La page ne dit pas
comment le produit est défini quand $2$ n'est pas inversible. Il l'est :
sur $M_1 \otimes M_2$, il existe une unique forme quadratique de forme
polaire $b_1 \otimes b_2$ qui vaut $2\,q_1(x)\,q_2(y)$ sur $x \otimes y$,
et elle est strictement non dégénérée avec $b_1$ et $b_2$. La vérification
est de nous. La page ne dit pas non plus comment les $\lambda^{i}$ sont
définies sur les modules quadratiques ; quand $2$ est inversible, ce sont
les puissances extérieures, avec la forme bilinéaire
$\det\bigl(b(x_i, y_j)\bigr)$.}. Si $S$ est connexe et que $2$ n'est pas
inversible, les rangs sont pairs : en un point de caractéristique
résiduelle $2$, la forme polaire est alternée ($b(x, x) = 2q(x) = 0$) et non
dégénérée\footnote{La page dit « $2$ non inversible dans
$\mathcal{O}_S$ » ; l'argument demande un point de caractéristique
résiduelle $2$, ce qui revient au même pour un topos ayant assez de
points.}. Si $2$ est inversible, $WG(S)$ est un anneau unitaire.

\emph{La question.} Même alors, Grothendieck ignore si $WG(S)$ est un
$\lambda$-anneau, c'est-à-dire si l'on a les identités universelles
\[
\lambda^{i}(xy) = P_i\bigl(\lambda^{1}x, \dots, \lambda^{i}x ;
\lambda^{1}y, \dots, \lambda^{i}y\bigr), \qquad
\lambda^{i}(\lambda^{j}x) = Q_{i,j}\bigl(\lambda^{1}x, \dots,
\lambda^{ij}x\bigr),
\]
à coefficients entiers\footnote{La page arrête les arguments de $P_i$ à
$\lambda^{i-1}$ ; les polynômes universels font intervenir $\lambda^{i}$,
comme le montre déjà $P_2 = x^{2}\lambda^{2}y + y^{2}\lambda^{2}x -
2\lambda^{2}x\,\lambda^{2}y$. On corrige.}. Il réduit la question à un
cas universel. Pour la première identité, il suffit de la vérifier sur le
topos classifiant $B_{O(m) \times O(n)}$, pour les images inverses des
modules quadratiques universels de rangs $m$ et $n$ : toute paire de modules
quadratiques de rangs $m$ et $n$ en provient par image inverse, et les deux
membres, biadditifs au sens des $\lambda$-anneaux, se déterminent sur les
classes effectives. De même la seconde se teste sur $B_{O(m)}$. « D'où la
tâche de déterminer $WG(B_{O(m)})$ », que le dossier n'entreprend
pas\footnote{Le « topos classifiant » s'entend ici comme celui qui classifie
les modules quadratiques de rang donné ; la page ne précise pas au-dessus
de quelle base. Sur le sort de la question : on sait aujourd'hui — nous le
citons de mémoire — que la réponse est positive pour un corps de
caractéristique différente de $2$ (S.~McGarraghy) et, plus généralement,
pour les schémas sur un tel corps (M.~Zibrowius, 2015), par une réduction
aux représentations des groupes orthogonaux proche de celle que la page
propose.}.

\pagerange{6}{7}
\subsection*{L'application hyperbolique et le groupe de Witt (pages 6 et 7)}

Soit $K_s(S)$ le groupe de Grothendieck « spécial » des modules localement
libres de type fini, construit sur les seules classes d'isomorphie, et
$K(S)$ son quotient par les relations issues des suites exactes courtes.
L'application $E \mapsto H(E)$ définit un homomorphisme
\[
\alpha \colon K_s(S) \longrightarrow WG(S).
\]
La page note qu'il n'est pas clair qu'il se factorise par $K(S)$,
c'est-à-dire que $H(E) = H(E') + H(E'')$ dans $WG(S)$ pour une extension
non scindée de $E''$ par $E'$. Il ne se factorise pas en
général\footnote{Contre-exemple de nous. Le rang et la classe du module
sous-jacent définissent un homomorphisme $WG(S) \to K_s(S)$, qui envoie
$H(E)$ sur $[E] + [\check{E}]$. Sur une courbe elliptique $X$, soit $F_2$
l'extension non scindée de $\mathcal{O}_X$ par $\mathcal{O}_X$ (Atiyah,
1957) ; elle est autoduale et indécomposable, et comme les fibrés sur une
variété complète vérifient le théorème de Krull–Schmidt (Atiyah, 1956),
$K_s(X)$ est libre sur les classes d'indécomposables. Donc
$[F_2] + [\check{F}_2] = 2[F_2] \neq 4[\mathcal{O}_X]$, et
$H(F_2) \neq 2H(\mathcal{O}_X)$ dans $WG(X)$. La définition moderne du
groupe de Grothendieck–Witt impose précisément cette relation.}.

Par définition, le \emph{groupe de Witt} $W(S)$ est le conoyau de $\alpha$ :
\[
K_s(S) \xrightarrow{\ \alpha\ } WG(S) \longrightarrow W(S) \longrightarrow 0 .
\]

\emph{Les formes $q \perp -q$.} Si $2$ est inversible, $F = (M, q) \perp
(M, -q)$ est hyperbolique, isomorphe à $H(M)$ : la diagonale $\Delta$ et
l'antidiagonale $\nabla$ de $M \oplus M$ sont totalement isotropes,
supplémentaires, et la forme polaire les met en dualité par
$\bigl((x, x), (y, -y)\bigr) \mapsto 2b(x, y)$. La page 7 ajoute que, sur un
schéma affine, « il paraît (Laborde dixit) » que c'est vrai sans condition
sur $2$, et que ce n'est pas vrai sur tout topos localement annelé, par
exemple sur $B_{O(n)}$, $n \geqslant 1$.

L'énoncé attribué à Laborde est juste, et se démontre ainsi. Supposons
qu'il existe une forme bilinéaire $\beta$ sur $M$, non nécessairement
symétrique, telle que $\beta(y, y) = q(y)$ ; c'est le cas si $M$ est
projectif, donc sur un schéma affine. La diagonale $\Delta$ est un
lagrangien sur lequel la forme quadratique $q \perp -q$ s'annule. Définissons
$g \colon M \to M$ par $b(x, g(y)) = -\beta(y, x)$ ; alors le sous-module
$\Lambda = \{(y + g(y), g(y))\}$ est un supplémentaire de $\Delta$, et
$(q \perp -q)(y + g(y), g(y)) = q(y) + b(y, g(y)) = 0$. Deux lagrangiens
supplémentaires sur lesquels la forme s'annule font de $F$ le plan
hyperbolique $H(\Delta) \simeq H(M)$\footnote{Démonstration de nous ; pour
$2$ inversible on prend $\beta = b/2$ et l'on retrouve l'argument de la page
6. L'existence de $\beta$ est une condition suffisante, dont nous ne
montrons pas qu'elle est nécessaire : l'affirmation de la page sur
$B_{O(n)}$, qui ne peut porter que sur une base où $2$ n'est pas inversible,
reste la sienne, sans démonstration.}.

\emph{Le groupe de Witt, directement.} Dans l'un et l'autre cas —
$2$ inversible, ou $S$ affine — toute classe de $W(S)$ est celle d'un module,
puisque $[M] - [M'] = [M \perp (M', -q')] - \alpha(M')$, et deux modules ont
même classe si et seulement s'il existe des modules localement libres $E$,
$F$ avec
\[
(M, q) \perp (M', -q') \perp H(F) \simeq H(E),
\]
ce que la page marque « à vérifier » en marge et qui se vérifie ainsi : une
égalité dans le groupe de Grothendieck est un isomorphisme après addition
d'un module $X$, et $X \perp (X, -q_X) \simeq H(X)$ se range dans les
termes hyperboliques. Sur un corps, où vaut le théorème de simplification
de Witt, la condition devient : $(M, q) \perp (M', -q') \simeq H(G)$ pour un
$G$\footnote{Le théorème de simplification de Witt vaut pour les formes
quadratiques non singulières sur un corps de toute caractéristique. La page
l'écrit « si on dispose d'un “th.\ de Witt” — ce qui est toujours le cas si
$S$ est le spectre d'un corps », « toujours » et « corps » étant des
lectures douteuses.}.

\emph{Structure de $W(S)$.} L'image $\alpha(K_s(S))$ est un idéal de
$WG(S)$, puisque $(N, q_N) \otimes H(E) \simeq H(N \otimes E)$ ; $W(S)$ est
donc un anneau\footnote{La page écrit $\alpha(K(S))$ ; il s'agit de
$\alpha(K_s(S))$.}. Mais il n'est plus augmenté vers $H^{0}(S,
\mathbb{Z})$, seulement vers $H^{0}(S, \mathbb{Z}/2\mathbb{Z})$, par le rang
modulo $2$, puisque les plans hyperboliques sont de rang pair. La page
demande si la pré-$\lambda$-structure passe au quotient. Elle n'y passe
pas : $\lambda^{2}\bigl(H(\mathcal{O}_S)\bigr) = \lambda^{2}\langle 1, -1
\rangle = \langle -1 \rangle$ est de rang $1$, donc non nul dans $W(S)$ dès
que $S$ est non vide et que $2$ y est inversible, alors que $H(\mathcal{O}_S)$
y est nul\footnote{Réponse de nous.}.

\pagerange{8}{8}
\subsection*{La classe de Stiefel–Whitney totale (page 8)}

On suppose maintenant l'anneau $\mathcal{O}_S$ strictement local (c'est le
cas de l'anneau structural d'un topos étale) et $2$ inversible\footnote{La
seconde hypothèse est ajoutée dans la marge ; « structural » et « ét. »
sont des lectures douteuses. Leur rôle, que la page ne dit pas : sur un
anneau local strictement hensélien où $2$ est inversible, toute unité est un
carré, de sorte que tout module quadratique non dégénéré de rang $n$ est
localement isomorphe à la forme unité $\langle 1, \dots, 1 \rangle$. Les
modules de rang $n$ sont alors les torseurs sous $O(n)$, et ceux de rang $1$
les torseurs sous $\mu_2 = \mathbb{Z}/2$.}. On définit alors la
\emph{classe de Stiefel–Whitney totale}
\[
w \colon WG(S) \longrightarrow 1 + \widehat{H}^{+}(S, \mathbb{F}_2),
\]
caractérisée, $S$ variant, par : a) sa fonctorialité en $S$ ; b) sa valeur
en rang $1$, $w(M, q) = 1 + \mathrm{cl}(M, q)$, où
$\mathrm{cl}(M, q) \in H^{1}(S, \mu_2) = H^{1}(S, \mathbb{F}_2)$ est la
classe du torseur $\mathrm{Isom}(\langle 1 \rangle, (M, q))$ ; c) la formule
$w(x + y) = w(x)\,w(y)$\footnote{La page annonce b) comme une
« normalisation » sans écrire la valeur ; celle-ci est la seule possible.
L'unicité repose sur un principe de scindage, que la page invoque plus bas
sous le nom d'« arguments habituels de splitting », et l'existence sur la
cohomologie des $B_{O(n)}$ ; ni l'une ni l'autre n'est établie ici. Sur un
corps, ces classes sont celles de A.~Delzant (1962) ; J.~Milnor s'en sert en
1970.}.

Un passage barré affirmait qu'« un calcul facile » donne $w \circ \alpha = 1$,
de sorte que $w$ se factoriserait par $W(S)$. C'est faux, et Grothendieck
l'a barré : $w\bigl(H(\mathcal{O}_S)\bigr) = w\langle 1, -1 \rangle = 1 +
[-1]$, et la classe $[-1]$ n'est pas nulle en général — sur le topos étale de
$\operatorname{Spec} \mathbb{R}$, par exemple. Le brouillon de la page 3
porte précisément sur $w(H(L))$, qu'il donne égal à
$1 + [-1] + c_1(L) \bmod 2$\footnote{Le
rapprochement avec la page 3 est de nous. La note barrée en marge commence
par « $w(\alpha(E)) = $ ».}.

Les arguments de scindage montrent en revanche que $w(M \otimes N)$ et
$w(\lambda^{i} M)$ s'expriment par des formules universelles en $w(M)$,
$w(N)$ et les rangs. Le rang et $w$ définissent donc un homomorphisme de
$\lambda$-anneaux
\[
WG(S) \longrightarrow H^{0}(S, \mathbb{Z}) \times \bigl(1 +
\widehat{H}^{+}(S, \mathbb{F}_2)\bigr),
\]
où le but porte la structure de $\lambda$-anneau définie par ces formules
universelles. C'est l'anneau $\widetilde{A}$ que le dossier 2 construit pour
les classes de Chern (pages 20 à 29), pris ici modulo $2$\footnote{Le
rapprochement avec le dossier 2 est de nous. La page écrit
$H^{0}(S, \mathbb{Z}) \times (1 + \hat{H}^{*+}(S, \mathbb{F}_2))$.}.

\pagerange{10}{16}
\subsection*{Les tirés à part (pages 10 à 16)}

Entre la page 8 et la suite, deux notes d'O.~Laborde aux \emph{Comptes rendus}
de 1974, sans annotation de sa main : « Théorème de Skolem-Nother pour les
algèbres d'Azumaya graduées sur un anneau semi-local » (pages 10 à 12) et
« Formes quadratiques, algèbres de Clifford et signatures » (pages 13 à 16).
La page 7 cite Laborde ; ces tirés à part n'étant pas transcrits, on ne
sait pas si l'énoncé qu'elle lui attribue s'y trouve. Ils sont la seule
date antérieure à 1975 que porte le dossier\footnote{Cette lecture ne les
résume pas.}.

\pagerange{17}{28}
\section*{III. Ensembles finis, groupes symétriques et opérations $\lambda$ (pages 17 à 28)}

Cette suite, au feutre noir puis à l'encre, n'a pas de titre. Elle part d'un
groupe $G$ et de ses ensembles finis, et aboutit aux représentations des
groupes symétriques.

\pagerange{17}{18}
\subsection*{Un brouillon : $\lambda$ et $\sigma$ sur les $G$-ensembles (pages 17 et 18)}

La page 17, un feuillet de listing, ne porte que $K_n(G) = K(G \times
\mathfrak{S}_n)$. La page 18 essaie sur un $G$-ensemble fini $E$ les
opérations des $\lambda$-anneaux : $\lambda^{i}(E)$ est la classe de
l'ensemble $\mathfrak{P}_i(E)$ des parties à $i$ éléments, $\sigma^{i}(E)$
celle de $\mathrm{Sym}^{i}(E) = E^{i}/\mathfrak{S}_i$, et la page encadre
$\sigma_{-t}(E)\,\lambda_t(E) = 1$, puis, en degré $2$,
\[
E \times E = \sigma^{2}(E) + \lambda^{2}(E), \qquad
\sigma^{2}(E) = \lambda^{2}(E) + E, \qquad
2\lambda^{2}(E) + E = E^{2},
\]
en vérifiant la dernière sur les cardinaux :
$2 \cdot \frac{d(d-1)}{2} + d = d^{2}$.

Ces identités sont vraies pour les cardinaux, c'est-à-dire pour $G$ trivial.
Dans l'anneau des $G$-ensembles virtuels, seule la deuxième subsiste, car
$\mathrm{Sym}^{2}(E)$ est la réunion de la diagonale et de
$\mathfrak{P}_2(E)$. Les deux autres tombent déjà pour $G = \mathbb{Z}/2$
agissant sur lui-même : $E \times E$ est formé de deux orbites libres, alors
que $2\lambda^{2}(E) + E$ est formé de deux points fixes et d'une orbite
libre ; et le coefficient de $t^{2}$ dans $\sigma_{-t}\lambda_t$ vaut
$2 - [G] \neq 0$\footnote{Calcul de nous. La raison est que
$E \times E \smallsetminus \Delta = \mathrm{Mon}(I_2, E)$, dont le quotient
par $\mathfrak{S}_2$ est $\mathfrak{P}_2(E)$, n'est pas en général
$2\,\mathfrak{P}_2(E)$ : le passage au quotient oublie l'action des
permutations.}. Ce qui se perd est l'action de $\mathfrak{S}_n$ sur les
$n$-uplets d'éléments distincts, et c'est précisément ce que les pages
suivantes conservent.

La page esquisse aussi le modèle linéaire, pour un élément $L$ « de rang
$1$ » : $\lambda_t(L) = 1 + tL$ et $\sigma_t(L) = (1 - tL)^{-1}$. Elle
encadre une flèche $\mathcal{E}(G) \hookrightarrow R(G) \to R_{\mathbb{Q}}(G)$
de l'anneau des $G$-ensembles vers l'anneau des représentations, qui,
pour $G$ fini, n'est injective que si $G$ est cyclique\footnote{Le rang de l'anneau de
Burnside est le nombre de classes de conjugaison de sous-groupes, celui de
$R_{\mathbb{Q}}(G)$ le nombre de classes de sous-groupes cycliques, et
l'application devient surjective après tensorisation par $\mathbb{Q}$
(théorème d'induction d'Artin). Elle est donc injective exactement quand
tout sous-groupe est cyclique, c'est-à-dire quand $G$ est cyclique.
L'indice du dernier $R$ est douteux.}. Elle se termine sur la question
« $\lambda^{i}(xy) = $ ? », sans réponse.

\pagerange{19}{20}
\subsection*{L'anneau $k_{*}(G)$ et les opérations $\tau^{n}$ (pages 19 et 20)}

Soit $G$ un groupe, $k_0(G)$ l'anneau des $G$-ensembles finis virtuels —
l'\emph{anneau de Burnside} de $G$ — et, pour $n \geqslant 0$,
\[
k_n(G) = k(G \times \mathfrak{S}_n), \qquad
k_{*}(G) = \bigoplus_{n \geqslant 0} k_n(G),
\]
avec $\mathfrak{S}_0 = \{1\}$ ; on note $\mathrm{cl}_n(E) \in k_n(G)$ la
classe d'un $G \times \mathfrak{S}_n$-ensemble fini. Le \emph{produit
d'induction}
\[
\mathrm{cl}_m(E) * \mathrm{cl}_n(F) = \mathrm{cl}_{m+n}\bigl((E \times F)
\times^{\mathfrak{S}_m \times \mathfrak{S}_n} \mathfrak{S}_{m+n}\bigr),
\]
où $G$ agit diagonalement, fait de $k_{*}(G)$ une $k_0(G)$-algèbre graduée,
commutative\footnote{La page écrit le produit contracté
$(E \times F) \wedge^{\mathfrak{S}_m \times \mathfrak{S}_n}
\mathfrak{S}_{m+n}$. La commutativité, que la page ne mentionne pas, vient
de la conjugaison par la permutation qui échange les deux blocs.}. Pour
$G$ trivial, c'est la somme des anneaux de Burnside des groupes symétriques,
munie du produit qui fait de $\bigoplus_n R(\mathfrak{S}_n)$ l'anneau des
fonctions symétriques.

Pour un $G$-ensemble fini $E$, on pose, avec $I_n = \{1, \dots, n\}$ et
$\mathfrak{S}_n$ agissant par la source,
\[
\tau^{n}(E) = \mathrm{cl}_n\bigl(\mathrm{Mon}(I_n, E)\bigr) \in k_n(G),
\qquad
\tau^{*}(E) = \sum_{n \geqslant 0} \tau^{n}(E) \in 1 +
\widehat{k_{*}(G)}^{+},
\]
où $\widehat{k_{*}(G)}^{+} = \prod_{n \geqslant 1} k_n(G)$ ; on a
$\tau^{0}(E) = 1$ et $\tau^{1}(E) = \mathrm{cl}_0(E)$, et la somme est finie
puisque $\mathrm{Mon}(I_n, E)$ est vide pour $n > \mathrm{card}\,E$.

\textbf{Proposition} (page 19). \emph{On a $\tau^{*}(E \sqcup F) =
\tau^{*}(E)\,\tau^{*}(F)$, c'est-à-dire $\tau^{n}(E \sqcup F) =
\sum_{i+j=n} \tau^{i}(E) * \tau^{j}(F)$.}

Une injection de $I_n$ dans $E \sqcup F$ est la donnée de la partie $J$ de
$I_n$ qui va dans $E$ et de deux injections $J \to E$,
$I_n \smallsetminus J \to F$ :
\[
\mathrm{Mon}(I_n, E \sqcup F) = \coprod_{i = 0}^{n}
\coprod_{J \in \mathfrak{P}_i(I_n)} \mathrm{Mon}(J, E) \times
\mathrm{Mon}(I_n \smallsetminus J, F),
\]
et chaque terme d'indice $i$ est stable par $\mathfrak{S}_n$. La page s'arrête
à « il suffit de voir que » la classe de ce terme est $\tau^{i}(E) *
\tau^{j}(F)$. C'est le cas : $\mathfrak{S}_n$ agit transitivement sur
$\mathfrak{P}_i(I_n)$, le stabilisateur de $I_i$ est $\mathfrak{S}_i \times
\mathfrak{S}_j$, et le terme est l'induit de $\mathrm{Mon}(I_i, E) \times
\mathrm{Mon}(I_j, F)$ de $\mathfrak{S}_i \times \mathfrak{S}_j$ à
$\mathfrak{S}_n$\footnote{Cette dernière étape est de nous.}.

\textbf{Corollaire} (page 20). \emph{Il existe un unique homomorphisme du
groupe additif $k_0(G)$ dans le groupe multiplicatif
$1 + \widehat{k_{*}(G)}^{+}$ tel que $\tau^{*}(\mathrm{cl}_0(E)) =
\tau^{*}(E)$ pour tout $G$-ensemble fini $E$.} C'est la propriété
universelle du groupe de Grothendieck, puisque $1 + \widehat{k_{*}(G)}^{+}$
est un groupe pour la multiplication.

La page définit enfin, pour un sous-groupe $H \subset \mathfrak{S}_n$,
l'application additive $q_H \colon k_n(G) \to k_0(G)$,
$\mathrm{cl}_n(E) \mapsto \mathrm{cl}_0(E^{H})$, et l'application
polynomiale de degré $n$
\[
q^{n}_H = q_H \circ \tau^{n} \colon k_0(G) \longrightarrow k_0(G).
\]
Le dossier n'en fait rien\footnote{La lecture ${}^{H}\!E$, les invariants
sous $H$, est incertaine. Les points fixes, comme les orbites, respectent
les réunions disjointes et définissent des applications additives ; avec
les orbites et $H = \mathfrak{S}_n$, on retrouverait
$\mathrm{Mon}(I_n, E)/\mathfrak{S}_n = \mathfrak{P}_n(E)$, le $\lambda^{n}$
de la page 18, alors qu'avec les points fixes $q_{\mathfrak{S}_n}\tau^{n}$
est nul pour $n \geqslant 2$. Les points fixes sont en revanche ceux qui
définissent les « marques » de Burnside.}.

\pagerange{21}{22}
\subsection*{Les puissances $T^{n}$, et ce qui les relie aux $\tau^{n}$ (pages 21 et 22)}

La page 21 définit de même, avec toutes les applications,
\[
T^{n}(E) = \mathrm{cl}_n\bigl(E^{I_n}\bigr) \in k_n(G), \qquad
T^{*}(E) = \sum_n T^{n}(E),
\]
multiplicatif sur les réunions disjointes par le même argument, et prolongé
à $k_0(G)$ comme au corollaire. C'est l'\emph{opération de puissance} de
l'anneau de Burnside, restreinte au sous-groupe diagonal $G \times
\mathfrak{S}_n$ du produit en couronne $G \wr \mathfrak{S}_n$\footnote{Les
opérations de puissance sont celles d'Atiyah en $K$-théorie (« Power
operations in $K$-theory », 1966), après la norme d'Evens (1963) en
cohomologie des groupes ; pour les anneaux de Burnside, elles sont étudiées
par tom Dieck (1979) et, sous le nom de transfert multiplicatif, par Tambara
(1993). La page ne cite personne.}. La page annonce une relation entre
$T^{*}$ et $\tau^{*}$. Elle part de la décomposition d'une application
$I_n \to E$ en la partition de $I_n$ par ses fibres et une injection du
quotient :
\[
E^{I_n} = \coprod_{Q \text{ quotient de } I_n} \mathrm{Mon}(Q, E) =
\coprod_{\alpha} \ \coprod_{Q \text{ de type } \alpha} \mathrm{Mon}(Q, E),
\]
où le \emph{type} $\alpha = (\alpha_1, \dots, \alpha_n)$ d'une partition
compte ses $\alpha_i$ classes à $i$ éléments, avec $\sum_i i\alpha_i = n$.

La formule encadrée qui suit, $T^{n}(E) = \sum_i \varphi_{n,i}(\tau^{i}(E))$
avec $\varphi_{n,i}(\mathrm{cl}_i(E)) = \mathrm{cl}_n(E
\times^{\mathfrak{S}_i} \mathfrak{S}_n)$, est annulée par deux grandes
croix, et à raison : l'induction de $\mathfrak{S}_i$ à $\mathfrak{S}_n$ ne
décrit pas l'action de $\mathfrak{S}_n$ sur les partitions. Voici la formule
juste\footnote{Elle est de nous. La page 22 écrit une somme indexée par les
types $\alpha$, $\sum_\alpha \varphi_{\nu(\alpha), \mu(\alpha)}\bigl(
\tau^{\mu(\alpha)}(x)\bigr)$ avec $\nu(\alpha) = n$ et $\mu(\alpha) = m$ ;
c'est la bonne indexation, mais l'opération $\varphi$ doit dépendre du type
et non des seuls $n$ et $m$.}. Soit $Q_\alpha$ une partition de type
$\alpha$ à $m = \sum_i \alpha_i$ classes. Son stabilisateur dans
$\mathfrak{S}_n$ est le produit de couronnes
$W_\alpha = \prod_i \mathfrak{S}_i \wr \mathfrak{S}_{\alpha_i}$, qui agit sur
l'ensemble des $m$ classes à travers
$\pi_\alpha \colon W_\alpha \to \prod_i \mathfrak{S}_{\alpha_i} \subset
\mathfrak{S}_m$. Alors
\[
T^{n}(E) = \sum_{\sum i\alpha_i = n}
\mathrm{Ind}_{G \times W_\alpha}^{G \times \mathfrak{S}_n}\,
\pi_\alpha^{*}\,
\mathrm{Res}^{G \times \mathfrak{S}_m}_{G \times \prod_i
\mathfrak{S}_{\alpha_i}}\, \tau^{m}(E).
\]
Le reste de la page 22 compte les partitions de type $\alpha$ :
\[
\frac{n!}{\prod_{i} \alpha_i!\,(i!)^{\alpha_i}},
\]
ce qui est juste. Sur les cardinaux, la formule redonne
$d^{n} = \sum_m S(n, m)\, d(d-1)\cdots(d-m+1)$, avec les nombres de Stirling
de seconde espèce. Un tableau de valeurs de $\tau^{2}(e_i)$ et
$\tau^{3}(e_i)$ est préparé et laissé vide.

\pagerange{23}{26}
\subsection*{Foncteurs sur la catégorie des ensembles finis et injections (pages 23 à 26)}

Soit $\mathcal{E}$ la catégorie des ensembles finis, les morphismes étant
les applications \emph{injectives}, et $\mathcal{H}$ celle des foncteurs
$\mathcal{E} \to (\mathrm{Ensf})$ à valeurs dans les ensembles finis. C'est
un prétopos\footnote{Lecture douteuse, mais l'énoncé est juste :
$\mathcal{H}$ est une sous-catégorie pleine du topos des foncteurs
$\mathcal{E} \to (\mathrm{Ens})$, stable par limites finies, sommes finies
et quotients par les relations d'équivalence, qui s'y calculent point par
point.}. On y distingue :
\begin{itemize}
  \item les foncteurs \emph{représentables} $h_I = \mathrm{Mon}(I, -)$ ;
  par Yoneda, $I \mapsto h_I$ identifie $\mathcal{E}^{\circ}$ à une
  sous-catégorie pleine de $\mathcal{H}$ ;
  \item les foncteurs $H_I = (-)^{I} = \mathrm{Hom}_{(\mathrm{Ens})}(I, -)$,
  restrictions à $\mathcal{E}$ des foncteurs correspondants sur
  $(\mathrm{Ensf})$.
\end{itemize}

\textbf{Proposition} (pages 23 et 24). \emph{Les $h_I$ sont des objets
connexes non vides de $\mathcal{H}$ ; à isomorphisme près, ce sont les
composantes connexes des $H_J$.} En effet $H_J = \coprod_Q h_Q$, la somme
portant sur les quotients $Q$ de $J$ : c'est, fonctorialisée, la
décomposition de la page 21\footnote{Le rapprochement est de nous. La
fonctorialité tient à ce qu'une injection $E \to E'$ ne change pas la
partition de $J$ par les fibres d'une application $J \to E$.}.

\emph{Foncteurs spéciaux} : les sommes finies de $h_I$ ; ce sont aussi les
facteurs directs des sommes finies de $H_I$\footnote{La page dit « des
sommandes directs (pour $I$ de cardinal assez grand) de foncteurs de la forme
$H_I$ ». Un seul $H_I$ ne suffit pas : $h_{I}$ apparaît dans $H_N$ avec pour
multiplicité le nombre de Stirling $S(\mathrm{card}\,N, \mathrm{card}\,I)$,
qui vaut $1$ si $\mathrm{card}\,I = 1$, de sorte que $2h_{\{1\}}$ n'est
facteur direct d'aucun $H_N$. On corrige en « sommes finies de ».}. Ils
sont stables par sommes, produits et plus généralement limites finies : un
produit $h_I \times h_J$ se décompose selon les configurations
d'images\footnote{$h_I \times h_J \simeq \coprod h_K$, sur les classes
d'isomorphie de couples d'injections $I \to K$, $J \to K$ conjointement
surjectives ; un égalisateur de deux flèches $h_I \rightrightarrows h_J$ est
$h_I$ ou vide. Ces vérifications sont de nous.}. Un foncteur spécial envoie
les injections sur des injections, donc s'identifie à un foncteur
$\mathcal{E} \to \mathcal{E}$, et le composé de deux foncteurs spéciaux est
spécial.

\emph{Foncteurs distingués} : les sommes finies de foncteurs
$E \mapsto \mathrm{Mon}(I, E)/H$, pour $H \subset \mathfrak{S}_I$ ; ce sont
aussi les sommes finies de composantes connexes des
$E \mapsto E^{I}/H$. La page affirme, sans le détailler, qu'ils ont les mêmes
propriétés de permanence que les foncteurs spéciaux. Une addition en marge
les décrit tous d'un coup :
\[
E \longmapsto \coprod_{n \geqslant 0} \mathrm{Mon}(I_n, E)
\times^{\mathfrak{S}_n} X_n,
\]
où $(X_n)$ est une suite de $\mathfrak{S}_n$-ensembles finis, presque tous
vides. C'est l'extension de Kan à gauche, le long de l'inclusion des
bijections dans les injections, de la suite $(X_n)$ — de ce qu'on appelle
depuis Joyal (1981) une \emph{espèce} finie. Dans le langage linéaire de
Church, Ellenberg et Farb (2015), les $\mathcal{E}$-ensembles sont des
« FI-ensembles », et les foncteurs distingués sont l'analogue ensembliste
des FI-modules induits, qu'on appelle aussi FI$\sharp$-modules\footnote{Les
deux noms sont postérieurs de plusieurs années à la page, qui ne pouvait pas
les connaître : les espèces de structures datent de 1981, les FI-modules de
2015. L'identification est de nous.}.

Un morphisme $u \colon F \to G$ de foncteurs $\mathcal{E} \to
(\mathrm{Ens})$ est une \emph{équivalence} s'il existe $n$ tel que $u(E)$ soit
bijectif pour tout $E$ de cardinal au moins $n$. Les équivalences vérifient
la propriété « deux sur trois » et sont stables par changement de base,
comme l'énonce la page ; ce sont les isomorphismes « en rang assez grand »,
le point de vue qui est aujourd'hui celui de la \emph{stabilité des
représentations} de Church et Farb\footnote{Le rapprochement est de
nous.}.

La page 26 pose enfin
\[
\Omega = \mathrm{End}(\mathcal{E}, \mathcal{E})
\subset \mathrm{Hom}(\mathcal{E}, (\mathrm{Ens})) =
\widehat{\mathcal{E}^{\circ}} ,
\]
la catégorie des foncteurs $\mathcal{E} \to \mathcal{E}$ — à valeurs finies,
envoyant les injections sur des injections — avec pour morphismes les
transformations naturelles à composantes injectives. C'est une
sous-catégorie non pleine, comme la page le souligne, du topos des
préfaisceaux sur $\mathcal{E}^{\circ}$. Elle porte trois opérations,
\[
(F + G)(E) = F(E) \sqcup G(E), \qquad (F \times G)(E) = F(E) \times G(E),
\qquad (F \circ G)(E) = F(G(E)),
\]
et le feuillet s'arrête là, le bas restant blanc. Une somme, un produit et
une composition qui se distribuent sont la structure d'un anneau de
composition, celle des opérations sur les $\lambda$-anneaux ; mais la page
ne passe pas au groupe de Grothendieck et ne dit pas où elle
va\footnote{Le rapprochement est de nous. Les anneaux de composition
remontent à Tall et Wraith (1970), qu'il pouvait connaître ; le nom de
« pléthorie » est de Borger et Wieland (2005).}.

\pagerange{27}{28}
\subsection*{Puissances tensorielles et représentations des groupes symétriques (pages 27 et 28)}

Soit $X$ un topos annelé commutatif, et
\[
K(X, *) = \bigoplus_{n \geqslant 0} K(B_{\mathfrak{S}_n X}),
\]
où $K(B_{\mathfrak{S}_n X})$ est le groupe de Grothendieck des modules
localement libres de type fini sur $X$ munis d'une action de
$\mathfrak{S}_n$, avec le produit d'induction. Pour $\xi = \mathrm{cl}(E)$
dans $K(X)$, on note $T^{n}(\xi)$ la classe de la puissance tensorielle
$E^{\otimes n}$, $\mathfrak{S}_n$ permutant les facteurs, et $T^{*}(\xi) =
\sum_n T^{n}(\xi) \in 1 + K(X, *)^{\wedge +}$\footnote{La page écrit
$\tau^{n}$ et $\tau^{*}$, et $\widetilde{\otimes}^{n}$ pour la puissance
tensorielle, lecture proposée. On écrit $T^{n}$ parce que c'est l'analogue
linéaire du $T^{n}$ de la page 21 : l'algèbre de $E^{I_n}$ est la
puissance tensorielle $n$-ième de l'algèbre de $E$.}. C'est le cadre du
dossier 2, pages 15 à 18, où le produit d'induction et les puissances
tensorielles d'un module sur un topos annelé sont déjà posés\footnote{Le
dossier 2 est daté par l'inventaire de 1957 ; rien ici ne le cite.}.

Soit $\rho_i$ l'élément unité de $K(B_{\mathfrak{S}_i}) = R_{\mathbb{Z}}
(\mathfrak{S}_i)$, regardé comme élément de degré $i$ de
$R(1, *) = K(e, *) = \bigoplus_i R(\mathfrak{S}_i)$, $e$ étant le topos
ponctuel ; le monôme $\rho^{\alpha} = \rho_1^{\alpha_1} \cdots
\rho_n^{\alpha_n}$ est la représentation de permutation de $\mathfrak{S}_n$
sur les classes du sous-groupe de Young $\prod_i \mathfrak{S}_i^{\alpha_i}$.
La page énonce :
\[
T^{n}(\xi) = \sum_{\sum i\alpha_i = n} m_\alpha\bigl(\lambda^{1}\xi, \dots,
\lambda^{n}\xi\bigr)\, \rho^{\alpha},
\]
où $m_\alpha \in \mathbb{Z}[S_1, \dots, S_n]$ est le polynôme isobare de
poids $n$ qui exprime la fonction symétrique monomiale
$\sum x_1^{\nu_1} \cdots x_m^{\nu_m}$ (avec $\alpha_i$ exposants égaux à
$i$) en les fonctions symétriques élémentaires $S_1, \dots,
S_n$\footnote{La page note ce polynôme $\Sigma_{\alpha_1, \dots, \alpha_n}$ ;
sa phrase s'interrompt au bas de la page 27 sur « en », et la page 28 ne la
reprend pas. On la complète.}.

La formule est juste, pour $E$ localement libre de type fini. Pour une
somme de fibrés en droites, $E^{\otimes n}$ se décompose selon les orbites de
$\mathfrak{S}_n$ sur les applications $I_n \to \{1, \dots, d\}$, et une
permutation des facteurs d'une puissance tensorielle d'une droite agit
trivialement ; le cas général s'en déduit par le principe de scindage, que
nous admettons pour $K(X, *)$. En degré $2$, elle donne $T^{2}(\xi) = \lambda^{2}\xi\,\rho_1^{2} + (\xi^{2} -
2\lambda^{2}\xi)\,\rho_2$, c'est-à-dire, si $2$ est inversible,
$\rho_1^{2}$ étant alors la représentation régulière $1 + \varepsilon$ de
$\mathfrak{S}_2$,
$T^{2}(\xi) = \sigma^{2}\xi \cdot 1 + \lambda^{2}\xi \cdot \varepsilon$ : la
suite exacte $0 \to \Lambda^{2}E \otimes \varepsilon \to E^{\otimes 2} \to
\mathrm{Sym}^{2}E \to 0$\footnote{Vérifications de nous.}. Par la
caractéristique de Frobenius, qui envoie $\rho_i$ sur la fonction
symétrique complète $h_i$, c'est l'identité de Cauchy
$\sum_\alpha m_\alpha(x)\,h_\alpha(y) = \prod_{i,j}(1 - x_i y_j)^{-1}$.

En marge, la page demande si, pour $\xi$ effectif, la formule est une somme
directe, et renvoie au théorème de commutation de H.~Weyl. Pas dans cette
base : $m_{(2)}(\xi) = \xi^{2} - 2\lambda^{2}\xi$ n'est pas effectif. La
somme directe existe dans l'autre base, celle des représentations
irréductibles : c'est la dualité de Schur–Weyl,
$E^{\otimes n} = \bigoplus_\mu S_\mu(E) \otimes V_\mu$, là où $n!$ est
inversible\footnote{Réponse de nous. H.~Weyl, \emph{The Classical Groups},
1939.}.

\textbf{Corollaire} (page 28). \emph{$\lambda^{n}(\xi)$ est le coefficient de
$\rho_1^{n}$ dans $T^{n}(\xi)$, et l'opération d'Adams $\psi^{n}(\xi)$ celui
de $\rho_n$.} En effet $m_{(1, \dots, 1)} = S_n$ et $m_{(n)} = \sum
x_j^{n}$ est la somme de Newton. C'est la démarche d'Atiyah (1966), qui tire
$\lambda^{n}$ et $\psi^{n}$ de l'opération de puissance\footnote{La page
écrit $S^{n}$ pour l'opération d'Adams. Elle ne cite pas Atiyah.}.

Un « coefficient » suppose que les $1_{K(X)} \cdot \rho^{\alpha}$ soient
linéairement indépendants sur $K(X)$, ce que la page croit « facile sans
doute ». C'est vrai si $X$ est au-dessus de $\mathbb{Q}$, et faux en
caractéristique $p$ : sur $\mathbb{F}_2$, la représentation régulière de
$\mathfrak{S}_2$ est extension de la triviale par la triviale, de sorte que
$\rho_1^{2} = 2\rho_2$\footnote{Contre-exemple de nous. Au-dessus de
$\mathbb{Q}$, les représentations de $\mathfrak{S}_n$ sont définies sur
$\mathbb{Q}$ et $K(B_{\mathfrak{S}_n X}) = K(X) \otimes
R_{\mathbb{Q}}(\mathfrak{S}_n)$, où les $\rho^{\alpha}$ forment une base.}.

\emph{Question} (page 28) : pour quels anneaux $k$ l'anneau
$R_k(1, *) = \bigoplus_{i \geqslant 0} R_k(\mathfrak{S}_i)$ est-il l'anneau
de polynômes en les $\rho_i$ ? Est-ce vrai pour un anneau local, ou pour
un anneau tel que $\mathbb{Z} \to K(k)$ soit un isomorphisme ? Pour un corps
de caractéristique nulle, oui : par Frobenius,
$\bigoplus_i R(\mathfrak{S}_i)$ est l'anneau des fonctions symétriques,
polynomial en les $h_i$. Pour un corps de caractéristique $p$, non : la
relation $\rho_1^{2} = 2\rho_2$ sur $\mathbb{F}_2$, et en général le rang de
$R_{\mathbb{F}_p}(\mathfrak{S}_n)$, qui est le nombre de partitions
$p$-régulières de $n$, l'interdisent. Un corps fini étant local, la réponse
à la seconde question est négative\footnote{Réponse de nous. La page
indexe la somme par « $i \in \mathbf{Z}$ » ; on lit $i \geqslant 0$.}.

\pagerange{29}{33}
\section*{IV. Le jeu de go (pages 29 à 33)}

Ces cinq pages, au feutre, n'ont de rapport avec rien d'autre dans le
dossier. Ce sont des mathématiques d'un jeu : une identité de comptabilité
entre deux façons de compter les points, et un essai de règles assez
précises pour qu'une partie finisse d'elle-même.

\pagerange{29}{29}
\subsection*{Deux décomptes, et le dernier coup (page 29)}

On joue une partie « complète » sur un goban de $\nu \times \nu$
intersections ; noir commence et joue $N$ coups, blanc en joue $N'$, avec
$\varepsilon = N - N' \in \{0, 1\}$, et $\varepsilon = 0$ si blanc joue le
dernier\footnote{La page écrit « $N' = N$ ou $N - 1$ », le signe étant
surchargé ; $N - 1$ est ce que demande le trait de noir.}. En fin de
partie, $P$ et $P'$ sont les nombres de pierres noires et blanches sur le
goban, $Q$ et $Q'$ les nombres de pierres noires et blanches capturées,
$T$ et $T'$ les territoires noir et blanc. Ainsi $N = P + Q$ et
$N' = P' + Q'$\footnote{La page appelle $Q$ le « nombre de prisonniers
noirs » ; l'égalité $N = P + Q$ impose qu'il s'agisse des pierres noires
prises par blanc, et non des pierres que noir a prises.}. On suppose que
toute intersection est, à la fin, une pierre ou un point de territoire :
\[
P + P' + T + T' = \nu^{2},
\]
ce qui exclut les points neutres et les seki, et demande que toutes les
pierres mortes aient été effectivement capturées — c'est le sens de
« complète »\footnote{Hypothèses implicites, rendues explicites ici. Un
seki est une position où des groupes des deux couleurs vivent ensemble sans
territoire ; ses libertés communes ne sont à personne.}.

Au décompte par territoire, noir marque $S = Q' + T$ et blanc
$S' = Q + T'$. Comme $Q' - Q = (N' - P') - (N - P) = (P - P') - \varepsilon$,
l'écart vaut
\[
D = S - S' = (P - P') + (T - T') - \varepsilon, \qquad \text{soit} \qquad
D = 2(P + T) - \nu^{2} - \varepsilon ,
\]
et symétriquement $D' = -D = 2(P' + T') - \nu^{2} + \varepsilon$\footnote{La
page écrit « $\nu^{2} - D' = 2(P' + T') + \varepsilon$ » ; c'est
$\nu^{2} + D'$, ou $\nu^{2} - D$, qui vaut $2(P' + T') + \varepsilon$. La
formule encadrée, elle, est juste.}. Or $(P + T) - (P' + T') = 2(P + T) -
\nu^{2}$ est l'écart au \emph{décompte par zone}, qui compte les pierres
posées et le territoire. L'identité dit donc que les deux décomptes, par
territoire (règles japonaises) et par zone (règles chinoises), diffèrent
exactement de $\varepsilon$ : d'un point si noir a joué le dernier, de rien
sinon\footnote{Le rapprochement avec les règles actuelles est de nous. Les
règles de l'AGA (1991) imposent, par des « pierres de passe », que blanc
joue le dernier, pour que les deux décomptes coïncident — nous le citons de
mémoire.}.

\pagerange{30}{30}
\subsection*{La position réglée (page 30)}

Supposons atteinte une position où les territoires des deux camps sont
« inattaquables », blanc ayant joué le dernier, de sorte que les deux camps
ont joué le même nombre $n = n'$ de coups et que noir a le trait. Avec les
mêmes lettres en minuscules — les prisonniers étant cette fois « virtuels »,
les pierres mortes restées sur le goban étant comptées comme prises — on
trouve de même, puisque $n = n'$,
\[
d = s - s' = 2(p + t) - \nu^{2} .
\]
Si chacun ne joue plus que dans son propre territoire, y compris pour y
capturer effectivement les pierres mortes, chaque coup ajoute une pierre et
ôte un point de territoire, et une capture rend au territoire des points
qu'il comptait déjà : $p + t = P + T$ est invariant, d'où
$D = d - \varepsilon$. La page ajoute que, faute de pouvoir prévoir
$\varepsilon$ — « est-elle indépendante de la suite du jeu ? » —, il est
raisonnable de compenser l'avantage du trait en posant $\varepsilon = 1$,
c'est-à-dire en comptant comme si noir jouait le dernier : une compensation
d'un point au décompte par zone. La compensation qu'on pratique
aujourd'hui, le komi, est de six à sept points\footnote{La page ne parle pas
de komi. « Inattaquables » est une lecture douteuse, et une ligne sur les
prisonniers virtuels est en partie illisible.}.

\pagerange{31}{32}
\subsection*{Règles de prise, et fin de l'obligation de jouer (pages 31 et 32)}

\emph{Prise.} Une pierre ne peut être posée sur une intersection libre où le
groupe qui la contiendrait n'aurait aucune liberté — le \emph{suicide} est
interdit —, sauf dans deux cas : le ko, renvoyé « plus bas » et jamais
décrit ; et le « remplissage d'un trou » : $A$ joue sur une intersection
$i$ qui est la dernière liberté d'un groupement $E$ de groupes de $B$, $E$
entourant le groupe de $A$ qui contient $i$. C'est la règle actuelle : un
coup qui prend est légal, les pierres prises étant retirées avant qu'on
examine les libertés de la pierre posée\footnote{Dans cette formulation, la
reprise d'un ko est elle-même un remplissage de trou ; ce que la règle du ko
ajoute est l'interdiction de reprendre aussitôt. La page n'en dit pas
plus.}.

\emph{Fin de partie.} Les joueurs jouent à tour de rôle et ne peuvent pas
passer. L'obligation de jouer ne tombe que dans un cas : le joueur $A$ n'a
plus d'intersection libre où il puisse légalement jouer, sinon des
intersections $c$ de la forme suivante : un groupement $E$ de $A$ n'a plus que
deux libertés isolées, toutes ses autres libertés étant prises par $B$, et
$c$ est l'une d'elles. Autrement dit, les seuls coups restants de $A$
bouchent l'un des deux yeux d'un de ses groupes. La page explique pourquoi
l'obligation doit alors tomber : si $A$ bouchait un œil, $B$ prendrait le
groupe en jouant dans l'autre, et, l'obligation persistant, $A$ perdrait de
proche en proche toutes ses pierres, « même si $A$ avait un avantage
considérable ».

Elle exige de plus que le groupement soit « inattaquable » : si c'était à
$B$ de jouer, il ne pourrait se placer ni en $c$ ni en $c'$, faute d'un
sous-groupement de $E$ entourant l'un d'eux et n'ayant plus que lui pour
liberté, et le ko ne le lui permettrait pas non plus. C'est la distinction
entre un vrai œil et un \emph{faux œil}\footnote{Le mot « isolées » est une
lecture douteuse. Le nom de faux œil est de nous.}.

\pagerange{33}{33}
\subsection*{Qui joue le dernier (page 33)}

En fin de partie, noir a $\rho$ groupements vivants, chacun avec au moins
deux yeux, et blanc $\rho'$. Partant de la position réglée de la page 30, où
noir a le trait, la page énonce : si $t - 2\rho \leqslant t' - 2\rho'$,
blanc joue le dernier ($\varepsilon = 0$, $D = d$) ; si
$t' - 2\rho' < t - 2\rho$, noir joue le dernier ($\varepsilon = 1$,
$D = d - 1$).

La règle est juste sous une hypothèse qu'elle ne dit pas : que chacun, à
partir de la position réglée, ne fasse que remplir son propre territoire
jusqu'à ne laisser que deux yeux d'un point à chaque groupe. Noir dispose
alors de $t - 2\rho$ coups et blanc de $t' - 2\rho'$ ; noir ayant le trait,
celui qui s'épuise le premier est noir si $t - 2\rho \leqslant t' - 2\rho'$,
et c'est alors blanc qui a joué le dernier\footnote{Vérification de nous.
L'obligation de jouer forcerait en fait le joueur qui s'épuise à jouer dans
le territoire adverse, là où une pierre a des libertés ; la page ne tient
pas compte de ces coups. Un passage barré affirmait que, blanc jouant le
dernier, tous ses groupes ont exactement deux yeux.}. La question de la page
30 reçoit ainsi une réponse partielle : $\varepsilon$ se lit sur la position
réglée, pourvu qu'on joue comme on vient de dire.

\pagerange{34}{39}
\section*{V. Groupes formels (pages 34 à 39)}

La page 34 est une chemise qui ne porte que ces mots ; la page 35 les
répète en tête. La suite est à l'encre bleue, et c'est la plus difficile à
lire du dossier : au-delà des formules, la page 35 est pour l'essentiel
une suite de lectures douteuses\footnote{On s'en tient ici à ce que la
transcription donne comme lisible, et l'on présente les énoncés comme les
siens, sans les démontrer.}. La page ne précise ni le corps de base, ni si
les groupes formels sont commutatifs ; les énoncés ne prennent leur sens qu'en
caractéristique $p > 0$, pour des groupes formels non nécessairement
commutatifs, où « $p$-divisible » désigne les groupes de Barsotti–Tate
connexes\footnote{Hypothèses supposées par nous. Les groupes de
Barsotti–Tate sont, sous l'angle galoisien, l'objet du dossier 161-5 (pages
22 et 23). Nous lisons « a-torique » comme « sans partie torique »,
c'est-à-dire sans partie de type multiplicatif ; la page ne le définit
pas.}.

\pagerange{35}{35}
\subsection*{Les énoncés de départ (page 35)}

La page énonce, sans démonstration :
\begin{itemize}
  \item (I) une extension d'un groupe $p$-divisible par un groupe lisse, ou
  unipotent, est triviale ; d'où, pour un $p$-divisible a-torique, un
  produit direct\footnote{« un produit direct » est une lecture douteuse.} ;
  \item tout sous-groupe $p$-divisible a-torique est central ; suit, en
  accolade, la liste de ce qu'il faudrait établir pour le plus grand
  sous-groupe $p$-divisible a-torique : son existence, sa stabilité par
  extension du corps de base, son comportement vis-à-vis des sous-groupes et
  des quotients ;
  \item (II) une extension d'un groupe formel lisse par un $p$-divisible est
  un produit « à isogénie près » ; d'où
  $D(G) \cap \mathrm{Cent}(G)$ quasi-unipotent, « primaire au $p$ », et
  mieux, $\mathrm{Cent}(D(G))$ quasi-unipotent.
\end{itemize}
La démonstration esquissée du corollaire de (II) prend le plus grand
sous-groupe $p$-divisible a-torique $T$ de $G$ : on aurait
$G \simeq T \times G/T$, donc $D(G) \cap T = \{e\}$. Un dernier corollaire
voudrait $G$ lisse isogène à un produit d'un facteur a-torique et d'un
autre facteur, et s'achève sur « … ? » ; une remarque, un NB et une note
en marge, qui disent qu'un certain quotient est « linéaire », ne se lisent
que par fragments\footnote{On ne les reconstruit pas : la transcription y
compte plus de mots illisibles que de mots lus.}. La page pose alors la
\emph{Question} : $D(G) \cap \mathrm{Cent}(G)$ est-il unipotent,
c'est-à-dire ne contient-il pas de $\mu_p$ ? Une note en marge ajoute qu'il
suffirait de le vérifier pour un quotient linéaire, et que
$\mathrm{Cent}(G/\mathrm{Cent}\,G)$ est « mieux » quasi-unipotent.

\pagerange{35}{36}
\subsection*{La « Solution » : un groupe vu par son centre et son groupe dérivé (pages 35 et 36)}

Sous le mot « Solution », la page introduit six groupes :
\[
L_1 = \mathrm{Cent}(G), \quad D_1 = D(G), \quad
\Pi_1 = L_1 \cap D_1, \quad \Pi_0 = G / L_1 D_1,
\]
et les deux quotients $G/\mathrm{Cent}(G)$ et $G/D(G)$. Ils forment deux
suites exactes, qui ont les mêmes extrémités :
\[
1 \to \Pi_1 \to \mathrm{Cent}(G) \longrightarrow G/D(G) \to \Pi_0 \to 1,
\qquad
1 \to \Pi_1 \to D(G) \longrightarrow G/\mathrm{Cent}(G) \to \Pi_0 \to 1 .
\]
La seconde flèche du milieu est un \emph{module croisé} : $G/\mathrm{Cent}(G)$
agit sur $D(G)$ par conjugaison, et c'est l'action $\theta$ que la page
dessine. La première l'est aussi, pour l'action triviale\footnote{La page
nomme $L_0 = G/\mathfrak{Z}(G)$ et $D_0 = G/D(G)$ et dessine les flèches
$L_1 \to L_0$ et $D_1 \to D_0$, avec $\theta \colon D_0 \to \mathrm{Aut}\,D_1$.
C'est incohérent : $\mathrm{Cent}(G) \to G/\mathrm{Cent}(G)$ est trivial, et
seule l'association $\mathrm{Cent}(G) \to G/D(G)$, $D(G) \to
G/\mathrm{Cent}(G)$ donne pour noyau $\Pi_1$ et pour conoyau $\Pi_0$ ; c'est
aussi celle que dessinent les deux arcs de la page 36, qui joignent $L_1$ à
$D_0$ et $D_1$ à $L_0$. On évite ici les lettres $L_0$, $D_0$.}. La
page note deux contraintes, qu'elle tire vraisemblablement des énoncés de
départ : $\Pi_1$ n'a pas de partie $p$-divisible, et $\Pi_0$ pas de
partie a-torique.

La « Solution » est donc un programme : reconstituer $G$ à partir des deux
modules croisés, qui partagent leurs invariants $\pi_1 = \Pi_1$ et
$\pi_0 = \Pi_0$. La page en donne la forme : les deux données se recollent
en
\[
D' = D(G) \sqcup_{\Pi_1} \mathrm{Cent}(G) = \mathrm{Cent}(G) \cdot D(G),
\qquad
L' = G/\mathrm{Cent}(G) \times_{\Pi_0} G/D(G) = G/\Pi_1,
\]
avec « obstruction nulle dans $H^{3}(\Pi_0, \Pi_1)$, indétermination dans
$H^{2}(\Pi_0, \Pi_1)$ »\footnote{Une note en marge identifie une image à
$\widehat{H}$, complété formel d'un groupe algébrique $H$ « canonique » :
c'est la situation du brouillon de la page 38.}. Elle se donne trois tâches : traiter le cas où
$\Pi_0$ est $p$-divisible ; dire quels couples de modules croisés sont
possibles ; donner des conditions pour que $G$ soit linéaire et $D(G)$
algébrisable. Seule la deuxième est reprise, aux pages 45 à 53, sous une
forme générale\footnote{Le rapprochement avec les pages 45 à 53 est de
nous ; il s'appuie sur l'identité des notations ($\varphi_G$, le diagramme
en losange, $H^{3}$ et $H^{2}$).}.

La page 36 amorce le calcul. Dans la situation générale
$1 \to \Pi_1 \to G_1 \to G_0 \to \Pi_0 \to 1$, avec $G_1 = \mathrm{Cent}(G)
\cdot D(G)$ et $G_0 = G/\Pi_1$, il y a un homomorphisme $\varphi_G$ de
$\mathrm{Cent}(G_0) \cap \operatorname{Ker}\theta$ vers
$Z^{1}(\Pi_0, \Pi_1)$ — ici vers $\mathrm{Hom}(\Pi_0, \Pi_1)$ —, de noyau
l'image de $\mathrm{Cent}(G)$ dans $G_0$. Dans le cas qui l'intéresse, la
page écrit cette application sur les éléments centraux de
$G/\mathrm{Cent}(G)$ qui agissent trivialement sur $D(G)$, à valeurs dans
$\mathrm{Hom}(\Pi_0, \Pi_1)$, et de noyau nul. D'où, $\Pi_1$ étant
quasi-unipotent, la conclusion encadrée : ce groupe d'éléments centraux est
quasi-unipotent. C'est, sous une forme encore floue, l'injectivité de
l'application $\Psi_G$ de la page 46, qui porte le même nom
$\varphi_G$\footnote{La phrase « ce qui exprime que $L_0 \hookrightarrow
\mathrm{Cent}\,G$ est iso » est une lecture douteuse. La page conclut
encore : « En fait, on doit avoir $G$ unipotent, comme $D_0 \simeq G/(G,G)$,
$\mathrm{Cent}(D_0)$ unipotent » ; nous ne savons pas la justifier et la
rapportons telle quelle.}.

\pagerange{37}{39}
\subsection*{Brouillons (pages 37 à 39)}

La page 37 aligne trois énoncés encadrés — $\mathrm{Cent}(D(G))$ unipotent,
$(G/\mathrm{Cent}\,G)_{\mathrm{ab}}$ linéaire, $\mathrm{Cent}(G/\mathrm{Cent}
\,G)$ unipotent — et une question entourée, « $(D(G))_{\mathrm{ab}}$
unipotent ? », avec une flèche vers « voir fin ». Aucune fin du dossier ne la
reprend.

La page 38 pose une extension centrale
\[
1 \to \mathfrak{z} \to D \to \widehat{H} \to 1,
\]
où $H$ est un groupe algébrique linéaire connexe lisse, $\widehat{H}$ son
complété formel et $\mathfrak{z}$ un groupe formel commutatif lisse : le
pendant formel de la structure que les pages 41 à 43 étudient pour les
groupes algébriques.

La page 39, couverte de schémas, calcule les extensions de
$\widehat{\mathbb{G}}_a$ par $\widehat{\mathbb{G}}_a$ par cocycles :
$\varphi(x, y)$, modifié par $\psi(x + y) - \psi(x) - \psi(y)$, dans
$H^{2}(\widehat{\mathbb{G}}_a, \widehat{\mathbb{G}}_a)$, et note, pour une
telle extension $E$, $\mathrm{Cent}(E) = D(E)$. Une telle extension non
commutative n'existe qu'en caractéristique $p$ : en caractéristique nulle,
un groupe formel est déterminé par son algèbre de Lie, et une extension
centrale de $k$ par $k$ est abélienne. En caractéristique $p$, la loi
$(x, z)(x', z') = (x + x', z + z' + x x'^{p})$ en donne une, de commutateur
$x x'^{p} - x' x^{p}$\footnote{Exemple de nous. La page ne dit pas quelle
extension elle considère.}.

\pagerange{41}{43}
\section*{VI. Groupes algébriques lisses connexes et décomposition de Chevalley (pages 41 à 43)}

Trois pages sans titre, à l'encre bleue, sur des groupes algébriques et non
formels. Aucune page ne dit à quoi elles servent ; elles sont rangées entre
les deux suites sur les groupes formels.

\pagerange{41}{41}
\subsection*{La structure (page 41)}

Soit $G$ un groupe algébrique lisse connexe sur un corps $k$, admettant une
décomposition de Chevalley
\[
1 \to L \to G \to A \to 1,
\]
$A$ variété abélienne, $L$ affine lisse connexe — c'est le cas si $k$ est
parfait (théorème de Chevalley)\footnote{On suppose désormais $k$ parfait,
ce que la décomposition $\mathfrak{z} = \mathfrak{z}_u \times \mathfrak{z}_m$
de la page 42 demande aussi. La page écrit « schéma abélien ».}. Soit
$M = G_{\mathrm{aff}}$ le plus grand quotient affine de $G$ et
\[
1 \to Z \to G \to M \to 1 .
\]
On sait que $Z$ est caractérisé par $M$ affine et $Z_{\mathrm{aff}} = 1$,
que $Z$ est central et que $Z \to A$ est surjectif, donc $G = Z \cdot L$ et
$L \to M$ surjectif\footnote{Ce sont des théorèmes de Rosenlicht (1956) ;
$Z$ est ce qu'on appelle aujourd'hui le sous-groupe \emph{anti-affine} de
$G$, étudié par M.~Brion (2009).}. Posons $\mathfrak{z} = L \cap Z$,
groupe affine commutatif. Alors $G \to M \times A$ est surjectif de noyau
$\mathfrak{z}$ central, et
\[
G = L \times^{\mathfrak{z}} Z = (L \times Z)/\mathfrak{z}
\]
est déterminé, avec le diagramme de ses sous-groupes, par les deux
extensions centrales : $L$ de $M$ par $\mathfrak{z}$, et $Z$ de $A$ par
$\mathfrak{z}$. La donnée de $G$ équivaut ainsi à celle de $M$, $A$,
$\mathfrak{z}$ et de deux extensions centrales, $L$ lisse connexe et
$Z_{\mathrm{aff}} = 1$ ; ou encore, en partant de $L$, à celle d'un groupe
affine lisse connexe $L$, d'une variété abélienne $A$, d'un sous-groupe
$\mathfrak{z} \subset \mathrm{Cent}(L)$ et d'une extension $Z$ de $A$ par
$\mathfrak{z}$ avec $Z_{\mathrm{aff}} = 1$\footnote{La page note ce $L$ avec
un astérisque, $L^{*}$ ; la reformulation n'en demande pas.}.

\pagerange{41}{42}
\subsection*{Extensions d'une variété abélienne par un groupe affine (pages 41 et 42)}

Une extension centrale de $A$ par $\mathfrak{z}$ commutatif est
commutative, puisque son commutateur définirait une application
bimultiplicative de $A \times A$ dans le groupe affine $\mathfrak{z}$, qui
est constante. Ces extensions sont données par la formule de Barsotti–Weil
généralisée :
\[
\mathrm{Ext}^{1}(A, \mathfrak{z}) =
\mathrm{Hom}\bigl(D(\mathfrak{z}), \mathrm{Pic}^{0}_{A/k}\bigr),
\]
où $D(\mathfrak{z})$ est le dual de Cartier de $\mathfrak{z}$ et
$\mathrm{Pic}^{0}_{A/k} = A^{*}$ la variété abélienne duale ; la catégorie de
ces extensions est rigide, puisque $\mathrm{Hom}(A, \mathfrak{z}) =
0$\footnote{La page écrit $\mathrm{Pic}_{A/k}$. C'est
$\mathrm{Pic}^{0}$ : pour $\mathfrak{z} = \mathbb{G}_m$, le fibré des
éléments non nuls d'un fibré en droites porte une loi de groupe au-dessus de
celle de $A$ si et seulement s'il est algébriquement équivalent à zéro
(Weil, Barsotti) ; voir le dossier 50, page 2. Pour la partie
infinitésimale de $D(\mathfrak{z})$, la page 42 note elle-même que les
homomorphismes se factorisent par $\mathrm{Pic}^{0}$. La commutativité des
extensions centrales est de nous.}. L'extension $Z$ est donc une flèche
$\varphi \colon D(\mathfrak{z}) \to A^{*}$, et il reste à exprimer
$Z_{\mathrm{aff}} = 1$.

\emph{$Z_{\mathrm{aff}}$ en termes de $\varphi$.} Le groupe $\mathfrak{z}$
se surjecte sur $Z_{\mathrm{aff}}$, puisque $A$ n'a pas de quotient affine
non trivial ; $Z_{\mathrm{aff}}$ est le plus grand quotient $\mathfrak{z}_1$
de $\mathfrak{z}$ tel que l'extension poussée de $A$ par $\mathfrak{z}_1$
soit triviale. Par dualité de Cartier, qui échange quotients de
$\mathfrak{z}$ et sous-groupes de $D(\mathfrak{z})$, c'est
$\mathfrak{z}_1 = D(\operatorname{Ker}\varphi)$. Donc
\[
Z_{\mathrm{aff}} = 1 \iff \varphi \text{ est un monomorphisme.}
\]
Décomposons $\mathfrak{z} = \mathfrak{z}_u \times \mathfrak{z}_m$ en ses
parties unipotente et de type multiplicatif, et soit $\Gamma =
D(\mathfrak{z}_m)$, groupe constant tordu de type fini ; $\varphi$ a deux
composantes $\varphi_0$ sur $D(\mathfrak{z}_u)$ et $\varphi_1$ sur $\Gamma$.
Leurs images, l'une connexe et l'autre étale, se coupent trivialement, de
sorte que $\varphi$ est un monomorphisme si et seulement si $\varphi_0$ et
$\varphi_1$ en sont.

\emph{Caractéristique nulle.} $\mathfrak{z}_u = W(t)$ est le groupe
vectoriel de son espace tangent $t$, et $\varphi_0$ est une application
linéaire $t^{\vee} \to H^{1}(A, \mathcal{O}_A)$, l'espace tangent de
$A^{*}$\footnote{La page écrit $\varphi_0 \colon t \to
t_{\mathrm{Pic}_{A/k}}$. Le dual de Cartier de $W(t)$ est le
groupe formel de $t^{\vee}$, et $\mathrm{Ext}^{1}(A, W(t)) = t \otimes
H^{1}(A, \mathcal{O}_A) = \mathrm{Hom}(t^{\vee}, H^{1}(A, \mathcal{O}_A))$ :
c'est $t^{\vee}$ qui se plonge.}. La donnée de $\mathfrak{z}$ et de $Z$ avec
$Z_{\mathrm{aff}} = 1$ équivaut donc à celle d'un sous-espace vectoriel de
$H^{1}(A, \mathcal{O}_A)$ et d'un sous-groupe constant tordu de type fini de
$A^{*}$.

\emph{Caractéristique $p$.} La page donne pour « bien connu » que la partie
unipotente de $Z_{\mathrm{aff}}$ est isogène à $\mathfrak{z}_u$, de sorte
que $Z_{\mathrm{aff}} = 1$ force $\mathfrak{z}_u$ fini ; son dual
$D(\mathfrak{z}_u)$ est alors fini radiciel, et $\varphi_0$ se factorise par
le complété formel $\widehat{A^{*}}$\footnote{L'énoncé « bien connu » est
cohérent avec le théorème de Brion (2009) selon lequel, en caractéristique
$p$, les groupes anti-affines sont des variétés semi-abéliennes ; nous le
citons de mémoire.}. La donnée de $\mathfrak{z}$ et de $Z$ équivaut à celle
d'un couple $(\tau, \Gamma)$ formé d'un sous-groupe radiciel $\tau$ et d'un
sous-groupe constant tordu $\Gamma$ de $A^{*}$, avec
$D(\mathfrak{z}) = \tau \times \Gamma \hookrightarrow A^{*}$.

\pagerange{43}{43}
\subsection*{Quels groupes $\mathfrak{z}$, et un défaut de constructibilité (page 43)}

\emph{NB.} Pour $k$ algébriquement clos, un groupe affine commutatif
$\mathfrak{z}$ intervient pour un $G$ convenable si et seulement s'il se
plonge centralement dans un groupe affine lisse connexe — ce qui, dit la
page, est toujours possible, et même dans un groupe commutatif — et si
$D(\mathfrak{z})$ se plonge dans une variété abélienne duale. D'après la
page, la seconde condition est toujours remplie en caractéristique nulle, et
l'est en caractéristique $p$ si et seulement si $\mathfrak{z}_u$ est fini. Il
faut ajouter, en caractéristique $p$, que $k$ ne soit pas une clôture
algébrique d'un corps fini : sur $\overline{\mathbb{F}}_p$, tout point d'une
variété abélienne est de torsion, et $\Gamma$ doit être fini, c'est-à-dire
$\mathfrak{z}_m$ fini\footnote{La restriction est de nous ; elle vient de
la correction $\mathrm{Pic}^{0}$ de la page 41. Elle s'accorde avec ce que
nous croyons être une remarque de Brion : sur $\overline{\mathbb{F}}_p$, les
groupes anti-affines sont des variétés abéliennes. Une note entourée, sous
le 2°, souligne que ce fait « distingue ici les caractéristiques » ; sa
lecture est douteuse.}.

\emph{NB.} La formation de $G_{\mathrm{aff}}$, pour un schéma en groupes
lisse à fibres connexes sur une base $S$, n'est pas « constructible » en
général. Exemple de la page : $G$ extension d'un schéma abélien $A$ par
$\mathbb{G}_a$, non triviale au point générique $y$, de caractéristique
nulle, et triviale en une spécialisation $s$ de caractéristique $p$ ; alors
$(G_y)_{\mathrm{aff}} = 1$ et $(G_s)_{\mathrm{aff}} \simeq \mathbb{G}_a$. Le
phénomène est plus fort que l'exemple : par l'étude de la page 42, une
extension de $A$ par $\mathbb{G}_a$ n'a jamais $G_{\mathrm{aff}} = 1$ en
caractéristique $p$, puisque $\mathfrak{z}_u = \mathbb{G}_a$ n'est pas fini.
Sur une base dominant $\operatorname{Spec}\mathbb{Z}$, le lieu où
$G_{\mathrm{aff}}$ est trivial est donc contenu dans la fibre de
caractéristique nulle, qui n'est pas constructible. En caractéristique
résiduelle fixée, la page affirme au contraire que la construction est
constructible\footnote{Le renforcement est de nous. La page écrit « groupe
affine lisse à fibres connexes », ce que l'exemple, qui n'est pas affine,
contredit ; nous lisons « groupe lisse à fibres connexes ».}.

\emph{Reformulation.} Il est « maladroit », dit la page, de se donner
$\mathfrak{z} \subset \mathrm{Cent}(L)$ et un monomorphisme
$D(\mathfrak{z}) \hookrightarrow A^{*}$ ; mieux vaut se donner $L$, $A$ et une
extension de $A$ par $\mathrm{Cent}(L)$ tout entier, c'est-à-dire un
homomorphisme $D(\mathrm{Cent}\,L) \to A^{*}$ : en poussant l'extension $Z$
le long de $\mathfrak{z} \subset \mathrm{Cent}(L)$, on retrouve
$G = L \times^{\mathrm{Cent}(L)} Z'$ sans condition d'injectivité. La
présentation précédente n'a d'avantage, conclut-il, que d'expliciter la
donnée sans faire intervenir des choses « à la Cantor » trop
« dégueulasses » — et la page s'arrête sur un tiret\footnote{« Cantor » est
une lecture douteuse. Il vise sans doute les sous-groupes arbitraires de
type fini d'un groupe de points comme $A^{*}(k)$, qu'aucune condition
naturelle ne décrit.}.

\pagerange{44}{53}
\section*{VII. Fourbis auxiliaires cohomologiques pour groupes formels (pages 44 à 53)}

La page 44 est une chemise qui porte ce titre. Les pages 45 à 53 sont
paginées par lui de 1 à 7, les pages 47 et 49 étant des versos
administratifs. Malgré le titre, rien n'y est propre aux groupes formels :
les groupes sont des faisceaux en groupes sur un topos $X$, et la
cohomologie celle du topos classifiant $B_{\pi_0}$ au-dessus de $X$ — pour
des groupes abstraits, la cohomologie des groupes. On y traite en général le
problème posé par la « Solution » de la page 35.

\pagerange{45}{45}
\subsection*{Le problème : deux modules croisés, un groupe (page 45)}

Un \emph{module croisé} $\mathcal{C} = (N, M, d, \Theta)$ est un
homomorphisme $d \colon N \to M$ avec une action $\Theta$ de $M$ sur $N$ telle
que $d$ soit équivariant pour l'action de $M$ sur lui-même par conjugaison
et que $d(n)$ agisse sur $N$ comme la conjugaison par $n$. Son noyau
$\pi_1$ est central dans $N$, son image est distinguée, et $M$ agit sur
$\pi_1$ à travers $\pi_0 = \operatorname{Coker} d$\footnote{La notion est de
J.~H.~C.~Whitehead (1949) ; la classification des modules croisés par
$(\pi_0, \pi_1)$ et une classe de $H^{3}(\pi_0, \pi_1)$ est de Mac Lane et
Whitehead (1950). Sous le nom de Gr-catégories, c'est l'objet de la thèse
de Hoàng Xuân Sính (1975), préparée sous la direction de Grothendieck ;
voir le dossier 135. La page ne nomme ni les modules croisés ni ces
travaux.}.

Si $G$ est un groupe et $N$, $N'$ deux sous-groupes distingués qui commutent
entre eux, on obtient deux modules croisés,
\[
\mathcal{C} = \bigl(N \to M = G/N'\bigr), \qquad
\mathcal{C}' = \bigl(N' \to M' = G/N\bigr),
\]
les actions étant induites par la conjugaison ; ils ont les mêmes
invariants, $\pi_1 = N \cap N'$ et $\pi_0 = G/NN'$, avec la même action de
$\pi_0$ sur $\pi_1$\footnote{L'hypothèse que $N$ et $N'$ commutent est
ajoutée ici : sans elle, $G/N'$ n'agit pas sur $N$. Elle est remplie dans le
cas qui intéresse la suite, où $N' = \mathrm{Cent}(G)$.}. Le problème de la
page est l'inverse : étant donnés deux modules croisés $\mathcal{C}$,
$\mathcal{C}'$ (données I et II) et des isomorphismes
$\pi_i(\mathcal{C}) \simeq \pi_i(\mathcal{C}')$ compatibles aux actions
(donnée III), trouver $G$.

Les données déterminent deux groupes : le sous-groupe $H = NN'$ de $G$ et le
quotient $K = G/\pi_1$,
\[
H = (N \times N')/\pi_1, \qquad K = M \times_{\pi_0} M',
\]
$\pi_1$ étant plongé antidiagonalement : un carré cocartésien et un carré
cartésien, que la page dessine. Les lignes
$(\alpha, \alpha') \cdot (\lambda, -\lambda) = (\alpha \cdot \lambda,
-\alpha' \cdot \lambda)$ vérifient que l'action de $K$ préserve
l'antidiagonale, ce qui est la compatibilité III\footnote{La page note $K =
G/\mathfrak{z}$, le « 3 » bouclé désignant ici $N \cap N' = \pi_1$.}. Et
$H \to K$ est un nouveau module croisé, d'invariants $\pi_0$, $\pi_1$.
Trouver $G$, c'est trouver une extension de $K$ par $\pi_1$ contenant $H$ de
façon compatible. La page énonce :
\[
\text{obstruction dans } H^{3}(B_{\pi_0}/X, \pi_1), \qquad
\text{indétermination dans } H^{2}(B_{\pi_0}/X, \pi_1) =
\mathrm{Extop}(\pi_0, \pi_1),
\]
c'est-à-dire : il existe une solution si et seulement si une classe de
degré $3$ s'annule, et les solutions forment alors un torseur sous le groupe
des classes d'extensions de $\pi_0$ par $\pi_1$ avec l'action
donnée\footnote{La page l'énonce sans démonstration. C'est la forme de la
théorie d'Eilenberg et Mac Lane (1947) des extensions à noyau non
abélien, obstruction dans $H^{3}$ et torseur sous $H^{2}$, et la classe
d'obstruction est l'invariant de Mac Lane–Whitehead du module croisé
$H \to K$ : cette identification est de nous. Le dossier 135 (pages 7 à 10)
retrouve l'obstruction de Mac Lane par les 2-gerbes relatives sur $B_G$ au-dessus de
$S$, dans la même notation.}.

La page isole ensuite les deux conditions qui intéressent la « Solution » :
$N'$ est central si et seulement si l'action $\Theta'$ de $M'$ sur $N'$ est
triviale, et $N$ contient $D(G)$ si et seulement si $M' = G/N$ est
commutatif. La première réduit $\mathcal{C}'$ à un homomorphisme
$d' \colon N' \to M'$ de noyau central, avec $N'$ commutatif et $d'(N')$
central dans $M'$, sans autre structure, et elle force l'action de
$\pi_0$ sur $\pi_1$ à être triviale — condition à vérifier sur
$\mathcal{C}$. La seconde rend $M'$, donc $\pi_0$, commutatif\footnote{En
marge de la page, quelques lignes en allemand, écrites verticalement, sans
rapport apparent avec la page ; non transcrites.}.

\pagerange{46}{46}
\subsection*{Que $N'$ soit exactement le centre (page 46)}

Supposons $N'$ central. L'image du centre de $G$ dans $M$ est contenue dans
le sous-groupe
\[
\mathfrak{Z} = \bigl[\mathrm{Cent}(K)/(N'/\pi_1)\bigr] \cap
\operatorname{Ker}\bigl(\Theta \colon M \to \mathrm{Aut}(N)\bigr),
\]
qui ne dépend que des données, et l'on a $N' \subset \mathrm{Cent}(G)
\subset \overline{N'}$, où $\overline{N'}$ est l'image inverse de
$\mathfrak{Z}$ dans $G$. Pour $g \in G$ et $z \in \overline{N'}$, le
commutateur $[g, z]$ tombe dans $\pi_1$ ; il est additif en $z$, additif en
$g$, nul pour $g \in N'$ et pour $g \in N$, puisque l'image de $z$ agit
trivialement sur $N$. Il définit donc un accouplement
$\pi_0 \times \mathfrak{Z} \to \pi_1$, c'est-à-dire
\[
\Psi_G \colon \mathfrak{Z} \longrightarrow \mathrm{Hom}(\pi_0, \pi_1),
\]
et $\mathrm{Cent}(G)$ est exactement l'image inverse de
$\operatorname{Ker}\Psi_G$. D'où :
\[
N' = \mathrm{Cent}(G) \iff \Psi_G \text{ est injectif.}
\]
La page passe par $\varphi_G \colon \pi_0 \to \mathrm{Hom}(\mathfrak{Z},
\pi_1)$, le même accouplement curryfié de l'autre côté\footnote{Les
vérifications d'additivité et de nullité sont de nous ; la page les
résume par « est triviale sur $N'$ et sur $\mathfrak{Z}$ ». En marge, une
définition antérieure, $\mathfrak{Z} = \operatorname{Ker}\Theta \cap
\mathrm{Cent}\,M$, est biffée.}.

\pagerange{48}{51}
\subsection*{Changer de solution : l'accouplement $c_E$ (pages 48 à 51)}

Une autre solution $G_1$ se déduit de $G$ par une extension $E$ de $\pi_0$
par $\pi_1$, ici centrale puisque l'action est triviale :
$G_1 = (G \times_{\pi_0} E)/\pi_1$. Le commutateur de deux relèvements dans
$E$ définit une forme alternée
\[
c_E \colon \pi_0 \times \pi_0 \longrightarrow \pi_1,
\]
bilinéaire parce que $\pi_0$ est commutatif, et nulle si et seulement si
$E$ est commutative\footnote{En marge : « on suppose $\pi_0$ commutatif ici.
Que dire sinon ? ». Dans le problème qui intéresse la « Solution », la
condition $N \supset D(G)$ le rend commutatif : la question ne s'y pose
pas.}. La page annonce, « au signe près »,
\[
\Psi_{G_1} = \Psi_G + \tilde{c}_E \circ \alpha,
\]
où $\alpha \colon \mathfrak{Z} \to M \to \pi_0$ est la projection et
$\tilde{c}_E \colon \pi_0 \to \mathrm{Hom}(\pi_0, \pi_1)$ la curryfication
de $c_E$. C'est juste : dans $G_1$, un commutateur est le produit des
commutateurs dans $G$ et dans $E$\footnote{Le « on doit trouver » de la
page est vérifié par nous. Le signe dépend du côté où l'on curryfie, ce que
la page note (« d'un signe ! »). Un $\beta$ biffé, qui envoyait
$\mathrm{Hom}(\pi_0, \pi_1)$ dans $\mathrm{Hom}(\pi_0, N')$, est abandonné
par la page elle-même.}. Il faut donc choisir $E$ de sorte que
$\Psi_G + \tilde{c}_E \alpha$ soit injectif.

\emph{Diviser la condition.} Soit $\mathfrak{Z}_0 = \operatorname{Ker}(\mathfrak{Z}
\to \pi_0)$, l'ensemble des éléments de $\mathfrak{Z}$ qui proviennent de
$H$. Sur $\mathfrak{Z}_0$, le terme correctif s'annule, et $\Psi_0 =
\Psi_G|_{\mathfrak{Z}_0}$ ne dépend ni de $G$ ni de $E$ : il se calcule dans
$H$, sur lequel $G$ agit à travers $K$\footnote{Une première tentative, par
une application $\varphi_0 \colon \mathfrak{Z} \to \mathrm{Hom}(\pi_0,
N'/\pi_1)$, est annulée de deux traits au bas de la page 48 et en tête de la
page 50 ; elle ne pouvait aboutir, $\Psi_G$ étant à valeurs dans
$\mathrm{Hom}(\pi_0, \pi_1)$. Une addition en marge identifie
$\mathfrak{Z}_0$ à $\mathrm{Cent}(N)/\pi_1 \cap \mathrm{Cent}(M)$, avec une
lecture douteuse.}. On obtient :
\begin{itemize}
  \item[b)] une condition préliminaire, indépendante de $E$ : $\Psi_0$ est
  injectif, c'est-à-dire $N' = H \cap \mathrm{Cent}(G)$, le sous-groupe des
  éléments de $H$ fixés par l'action de $K$\footnote{La page écrit « $N'$
  est égal au centre de $H$, $N' = \mathrm{Cent}(H)$ », avec un point
  d'interrogation en marge. Le noyau de $\Psi_0$ est
  $(H \cap \mathrm{Cent}\,G)/N'$ ; comme $\mathrm{Cent}(H)$ contient
  $H \cap \mathrm{Cent}(G)$, l'égalité $N' = \mathrm{Cent}(H)$ entraîne b),
  mais la réciproque n'a pas de raison d'être vraie. On corrige.} ;
  \item[c)] une condition qui dépend de $E$ : l'application
  $\overline{\Psi}_{G_1} \colon \mathfrak{Z}/\mathfrak{Z}_0 = \pi_0^{*} \to
  \mathrm{Hom}(\pi_0, \pi_1)/\Psi_0(\mathfrak{Z}_0)$ déduite de
  $\Psi_{G_1}$ est injective, $\pi_0^{*}$ désignant l'image de
  $\mathfrak{Z}$ dans $\pi_0$.
\end{itemize}
Les deux ensemble équivalent à l'injectivité de $\Psi_{G_1}$. La phrase qui
devait dire ce que c) impose à la classe de $E$ reste
inachevée\footnote{Elle s'arrête sur « Cette condition sur $G_1$ i.e.\ sur
$E_G \in \operatorname{Ob}\mathrm{Extop}(\pi_0, \pi_1)$ », avec
« (condition ?) » en marge.}.

\pagerange{51}{53}
\subsection*{Que $N$ soit exactement le groupe dérivé (pages 51 à 53)}

On veut maintenant $N = D(G)$. D'abord a) $N \supset D(G)$, c'est-à-dire
$M'$ commutatif, d'où $\pi_0$ commutatif. Soit
$\mathcal{D} = (N_{\mathrm{ab}})_M$ le groupe des coïnvariants sous $M$ de
l'abélianisé de $N$, et $N_0$ le noyau de $N \to \mathcal{D}$ ; tout quotient
$N/N_1$ sur lequel $G$ agit trivialement est un quotient de $\mathcal{D}$.
Dans $G/N_0$, extension centrale de $M'$ par $\mathcal{D}$, le commutateur de
deux relèvements définit une forme bilinéaire alternée, qui s'annule dès
qu'un des arguments est dans $d'(N')$ puisque $N'$ est central :
\[
\lambda_G \colon \pi_0 \otimes \pi_0 \longrightarrow \mathcal{D},
\qquad
N/D(G) \simeq \mathcal{D}/\lambda_G(\pi_0 \otimes \pi_0) .
\]
La condition cherchée est donc que $\lambda_G$ soit
surjectif\footnote{La page écrit « si $\bar{x}$ ou $\bar{y} \in d'(N')$ » ;
ce sont $x$ et $y$, non leurs relèvements, qui doivent être dans $d'(N')$.
L'isomorphisme $N/D(G) \simeq \mathcal{D}/\lambda_G(\pi_0 \otimes \pi_0)$
vient de ce que $D(G)/N_0$ est le groupe dérivé de $G/N_0$, engendré par ces
commutateurs.}. Quand on change $G$ en $G_1$ par l'extension $E$,
\[
\lambda_{G_1} = \lambda_G + \beta \circ c_E,
\]
où $\beta \colon \pi_1 \to \mathcal{D}$ est le passage au quotient.

La page divise encore. Soit $\mathcal{D}_0 = \operatorname{Coker}(\pi_1 \to
\mathcal{D})$ et $\lambda_0 \colon \pi_0 \otimes \pi_0 \to \mathcal{D}_0$
déduit de $\lambda_G$ : il ne dépend pas de $G$, puisqu'il se construit dans
$K = G/\pi_1$. Soit $P = \operatorname{Ker}\lambda_0$ et
$\pi_{1*} = \operatorname{Im}(\pi_1 \to \mathcal{D})$ ; $\lambda_G$ induit
$\overline{\lambda}_G \colon P \to \pi_{1*}$. Alors $\lambda_G$ est
surjectif si et seulement si
\begin{itemize}
  \item[b)] $\lambda_0$ est surjectif — condition indépendante de $G$, qui
  signifie que $N/\pi_1 = D(K)$, ou encore que $d(N) = D(M)$ ;
  \item[c)] $\overline{\lambda}_G \colon P \to \pi_{1*}$ est surjectif —
  condition qui dépend de $G$, par la formule
  $\overline{\lambda}_{G_1} = \overline{\lambda}_G + p\,c_E\,i$, où $i$ est
  l'inclusion $P \subset \pi_0 \otimes \pi_0$ et $p$ la projection
  $\pi_1 \to \pi_{1*}$.
\end{itemize}
L'équivalence de $N/\pi_1 = D(K)$ et de $d(N) = D(M)$, que la page marque
« vérifier », est exacte\footnote{Vérification de nous :
$d(N) = NN'/N'$ et $D(M) = D(G)N'/N'$ ; si $NN' = D(G)N'$, alors, $D(G)$ étant
contenu dans $N$, la loi modulaire donne $N = D(G)(N \cap N') =
D(G)\,\pi_1$, qui est $N/\pi_1 = D(K)$. La réciproque est immédiate. De
même, b) et c) ensemble équivalent à la surjectivité de $\lambda_G$.}.

Le dossier s'arrête là. Ce qu'il a obtenu tient en peu de mots : une fois
levée l'obstruction de la page 45, le groupe cherché est déterminé par le
choix d'une extension centrale $E$ de $\pi_0$ par $\pi_1$, et les deux
conditions c) ne dépendent de $E$ qu'à travers sa forme alternée $c_E$, de
façon affine. Pour des groupes abstraits, toute forme bilinéaire alternée
$\pi_0 \times \pi_0 \to \pi_1$ est la forme d'une extension centrale, et le
problème devient un problème d'algèbre linéaire : trouver une forme alternée
$c$ qui rende $\overline{\Psi}_G + \overline{\tilde{c}\alpha}$ injectif et
$\overline{\lambda}_G + p\,c\,i$ surjectif\footnote{Cette mise en forme,
et la surjectivité de $H^{2}(\pi_0, \pi_1) \to \mathrm{Alt}(\pi_0, \pi_1)$
pour des groupes abéliens abstraits, sont de nous. Pour des groupes formels
ou des faisceaux, la surjectivité demande une vérification que le dossier
ne fait pas.}. Le dossier ne le résout pas, et ne revient pas aux groupes
formels.

\end{document}
